% ===========================================================================
% Chapter 4: Generalizations (formalization plan)
% ===========================================================================
% Gacs source lines 4871-7289, reproduced in full with the author's permission.
% Nothing of this chapter is formalized yet, so no node carries \lean or \leanok.
% Planned nodes carry neither \notready nor \leanok; out-of-scope nodes carry
% \notready.  Our own infrastructure (Definitions A-C, Theorems D-F) uses the
% separate "ourdefinition" counter, so every number of his is unchanged.
%
% Changes to his text: sloppypar/sloppyenv wrappers dropped; one \qedhere
% dropped (proof of Proposition 4.6.10); references to Appendices A and B
% replaced by their printed numbers; numbered items he left unlabeled get
% labels <kind>:g<number>.

\chapter{Generalizations}

\begin{formalization}
Nothing in this chapter is formalized yet: this part of the blueprint is a plan. G\'acs works
over computable complete metric spaces, with arbitrary and possibly noncomputable measures, and
leans on the computable analysis of Appendix B. Neither Mathlib nor this repository has that
theory, so the plan specializes every statement to two spaces: the discrete space $\dS$ of finite
binary strings (identified with $\dN$), and the Cantor space $\dB^{\dN}$ of infinite binary
sequences, where Chapters 2 and 3 already work, cylinder by cylinder. A measure is a function on
strings, and a noncomputable measure reaches a computation through its names, which are oracles
in the sense of Chapter 3 (Definition~\ref{our:measure}). Lower semicomputability in a measure
argument is defined without reference to names (Definition~\ref{our:lsc}), and so is
$K(x \mvert \mu)$ (Definition~\ref{our:m-cond}). These and the facts from Appendix B that the
proofs use are stated as our Definitions A--C and Theorems D--F, which have their own counter.

Planned: Sections 4.1, 4.2, 4.4 and 4.5, Subsections 4.6.1, 4.6.3 and 4.6.4, and from
Section 4.3 Lemma~\ref{lem:concentr} and Proposition~\ref{thm:no-neutral}. Each planned node
lists the results of Chapters 1--3 it generalizes. Out of scope, marked not ready: the existence
of neutral measures and the semicomputability results of Section 4.3 (compactness of the space of
measures, Sperner's lemma), Proposition~\ref{thm:test-charac} on general metric spaces, and all of
Section 4.7.

Conventions as in Chapters 2 and 3: tests are multiplicative ($\t = 2^{\d}$, $2^{-H} = \t$),
valued in $\clint{0}{\infty}$, with no logarithms; constants are existential unless the statement
is about the constant, as $K(k)$ in Proposition~\ref{propo:test-computable-mon}.
\end{formalization}

 \section{Continuous spaces, noncomputable measures}\label{sec:cont-spaces}

This section starts the consideration of randomness in continuous spaces and
randomness with respect to noncomputable measures.

\subsection{Introduction}

The algorithmic theory of randomness is well developed when the underlying
space is the set of finite or infinite sequences and the underlying
probability distribution is the uniform distribution or a computable
distribution.
These restrictions seem artificial.  
Some progress has been
made to extend the theory to arbitrary Bernoulli distributions by
Martin-L\"of in~\cite{MLof66art}, and to arbitrary distributions, by Levin
in~\cite{LevinRand73,LevinUnif76,LevinRandCons84}.
The paper~\cite{HertlingWeihrauchRand98} by Hertling and Weihrauch 
also works in general spaces, but it is restricted to computable
measures.
Similarly, Asarin's thesis~\cite{Asarin88} defines randomness for sample
paths of the Brownian motion: a fixed random process with computable
distribution.

The exposition here has been inspired mainly by Levin's early 
paper~\cite{LevinUnif76} (and the much more elaborate~\cite{LevinRandCons84}
that uses different definitions): let us summarize part of the content
of~\cite{LevinUnif76}.
The notion of a constructive topological space \( \bX \) and the space of
measures over \( \bX \) is introduced.
Then the paper defines the notion of a uniform test.
Each test is a lower semicomputable function \( (\mu,x) \mapsto f_{\mu}(x) \),
satisfying \( \int f_{\mu}(x) \mu(dx) \le 1 \) for each measure \( \mu \).
There are also some additional conditions.
The main claims are the following.
 \begin{alphenum}

  \item There is a universal test \( \t_{\mu}(x) \), a test 
such that for each other test \( f \) there is a constant \( c > 0 \) with
\( f_{\mu}(x) \le c\cdot \t_{\mu}(x) \).
  The \df{deficiency of randomness} is defined as 
\( \d_{\mu}(x) = \log\t_{\mu}(x) \). 

  \item The universal test has some strong properties of ``randomness
conservation'': these say, essentially,
that a computable mapping or a computable
randomized transition does not decrease randomness.

  \item There is a measure \( M \) with the property that for
every outcome \( x \) we have \( \t_{M}(x) \le 1 \).
In the present work, we will call such measures \df{neutral}.

  \item\label{i:Levin.semimeasure}
Semimeasures (semi-additive measures)
are introduced and it is shown that there is a lower semicomputable
semi-measure that is neutral (let us assume that the \( M \) introduced above
is lower semicomputable).

  \item Mutual information \( I(x : y) \)
is defined with the help of (an appropriate
version of) Kolmogorov complexity, between outcomes \( x \) and \( y \).
It is shown that \( I(x : y) \) is essentially equal to \( \d_{M \times M}(x,y) \).
This interprets mutual information as a kind of ``deficiency of independence''.

 \end{alphenum}
This impressive theory leaves a number of issues unresolved:
  \begin{enumerate}

   \item The space of outcomes is restricted to be a compact topological
space, moreover, a particular compact space: the set of 
sequences over a finite alphabet (or, implicitly in~\cite{LevinRandCons84},
a compactified infinite alphabet).
However, a good deal of modern probability theory happens over
spaces that are not even locally compact:
for example, in case of the Brownian motion, over the
space of continuous functions.

   \item The definition of a uniform randomness test includes some conditions
(different ones in~\cite{LevinUnif76} and 
in~\cite{LevinRandCons84}) that seem somewhat arbitrary.

   \item No simple expression is known for the general
universal test in terms of description complexity.
Such expressions are nice to have if they are available.

  \end{enumerate}

Here, we intend to carry out as much of Levin's program as seems
possible after removing the restrictions.
A number of questions remain open, but we feel that they are worth
to be at least formulated.
A fairly large part of the exposition
is devoted to the necessary conceptual machinery.
This will also allow to carry further some other initiatives started in the
works~\cite{MLof66art} and~\cite{LevinRand73}: the study of tests that test
nonrandomness with respect to a whole class of measures (like the Bernoulli
measures).

Constructive analysis has been developed by several authors,
converging approximately on the same concepts,
We will make use of a simplified  version of the theory introduced  
in~\cite{WeihrauchComputAnal00}.
As we have not found a constructive measure theory in the
literature fitting our purposes, we will develop this theory here,
over (constructive) complete separable metric spaces.
This generality is well supported by standard results in measure
theoretical probability, and is
sufficient for constructivizing a large part of current probability
theory.

The appendix recalls some of the needed topology, measure theory,
constructive analysis and constructive measure theory.
We also make use of the notation introduced there.


\subsection{Uniform tests}

We first define tests of randomness with respect to an arbitrary measure.
Recall the definition of lower semicomputable
real functions on a computable metric space \( \bX \).

\begin{formalization}
Definitions A--C and Theorems D--E below are ours. They restate, for the discrete space $\dS$
and the Cantor space $\dB^{\dN}$, what G\'acs takes from Appendix B: the space of measures, lower
semicomputability on it, and conditional tests; Theorems D and E collect the facts from
Appendix B that his proofs use. Theorem F, after Remark~\ref{rem:oracle-warning}, relates
$K(x \mvert \mu)$ to Chapter 3's oracle complexity.
\end{formalization}

% Ours: the space of measures over our two spaces, and names (Appendix B, constructive measure theory)
\begin{ourdefinition}[Measures and their names]\label{our:measure}
Let \( \bX \) be either the discrete space \( \dS \) of finite binary strings or the Cantor
space \( \dB^{\dN} \) of infinite binary sequences.
 \begin{alphenum}
  \item A \df{measure} over \( \dS \) is a function \( \mu:\dS\to\clint{0}{1} \) with
\( \sum_{x}\mu(x) = 1 \).
A \df{measure} over \( \dB^{\dN} \) is a function \( \mu:\dS\to\clint{0}{1} \), giving the
measures of the cylinders, with \( \mu(\Lg) = 1 \) and \( \mu(w) = \mu(w0)+\mu(w1) \) for all
\( w \).
In both cases \( \cM(\bX) \) is the set of all such measures; no computability is assumed.
  \item A \df{name} of a measure \( \mu \) is an infinite binary sequence \( A \) that lists, in a
fixed order and a self-delimiting code, a rational number \( q_{A}(w,k) \) for every string
\( w \) and every \( k\in\dN \), with \( \card{\mu(w)-q_{A}(w,k)} \le 2^{-k} \).
Every finite prefix of \( A \) thus determines finitely many of the values \( q_{A}(w,k) \).
  \item The integral of \( f:\dS\to\clint{0}{\infty} \) is \( \mu f = \sum_{x}\mu(x)f(x) \).
Over \( \dB^{\dN} \), if \( f(\xi) = \sup_{n}g(\xi^{\le n}) \) for a function \( g \) on strings
that is monotonic with respect to the prefix relation, then
\( \mu f = \sup_{n}\sum_{w\in\dB^{n}}\mu(w)g(w) \).
 \end{alphenum}
\end{ourdefinition}

\begin{formalization}
Planned. These are Chapter 2's \texttt{ComputableMeasure} and \texttt{ComputableSeqMeasure}
(with \texttt{IsSeqMeasure}) without their computability field, so Chapter 2's computable
measures are a special case. A name is a \texttt{Seq}, read through its finite prefixes
\texttt{seqTake A m}, as Chapter 3's oracle machines read an oracle. Every measure has a name,
and a measure is computable in Chapter 2's sense (\texttt{IsComputableENNReal}) if and only if it
has a computable name. The integral over $\dB^{\dN}$ is Chapter 2's \texttt{levelIntegral} with
the supremum over $n$; by monotone convergence it does not depend on $g$. The counting measure
$\#$ on $\dS$, $\#(x) = 1$, serves in Section 4.5 as an integrating measure; it is not an element
of $\cM(\dS)$.
\end{formalization}

% Ours: lower semicomputability with a measure argument (Appendix B, semicomputability)
\begin{ourdefinition}[Lower semicomputability in a measure argument]\label{our:lsc}
  \uses{our:measure, def:semicomputability}
Let \( \bX \) be as in Definition~\ref{our:measure}.
 \begin{alphenum}
  \item A \df{basic set} is a subset of \( \bX\times\dS\times\cM(\bX) \) of the form
 \[
  \setof{(x,y,\mu) : x\in w,\; y = v,\; a_{j} < \mu(u_{j}) < b_{j} \txt{ for } j=1,\dots,k},
 \]
given by strings \( w,v,u_{1},\dots,u_{k} \) and rationals \( a_{j}<b_{j} \).
Here \( x\in w \) means \( x = w \) over \( \dS \) and \( w\prefix x \) over \( \dB^{\dN} \).
Basic sets are coded by natural numbers.
  \item A function \( f:\bX\times\dS\times\cM(\bX)\to\clint{-\infty}{\infty} \) is \df{lower
semicomputable} if there is a computably enumerable set \( E \) of pairs (basic set \( S \),
rational \( r \)) such that for every rational \( r \), the set \( \setof{f>r} \) is the union of
the basic sets \( S \) with \( (S,r)\in E \).
For \( f\ge 0 \) this says that \( f \) is the supremum of an enumerated family of functions
\( r\cdot 1_{S} \).
It is \df{upper semicomputable} if \( -f \) is lower semicomputable.
  \item For a fixed measure \( \mu \), a function \( g \) over \( \bX \) is
\df{\( \mu \)-lower semicomputable} if \( g = f(\cdot,\mu) \) for some lower semicomputable
\( f \) without the argument \( y \).
  \item The same definitions apply with the argument \( y \) absent or ranging over \( \dZ \),
and with several measure arguments, possibly over different spaces; each measure argument
contributes finitely many conditions to a basic set.
 \end{alphenum}
\end{ourdefinition}

\begin{formalization}
Planned. The definition sees a measure only through finitely many of its values, so it does not
depend on names; Remark~\ref{rem:oracle-warning} explains why this matters. Without a measure
argument it is Chapter 1's lower semicomputability (\texttt{IsLSC}, and Chapter 2's unary
version), and substituting a computable measure leads back there
(Theorem~\ref{our:lsc-props}). Over both spaces the sets $\setof{\mu : a<\mu(u)<b}$ generate the
weak topology on probability measures, so this is G\'acs's notion restricted to our spaces. In
Lean: a computable enumeration of codes of basic sets with rational levels, a decidable check of
finitely many strict inequalities, and the supremum over the enumeration.
\end{formalization}

% Ours: conditional universal tests, m(x | mu), K(x | mu)
\begin{ourdefinition}[Conditional tests; \( \m(x \mvert \mu) \) and \( K(x \mvert \mu) \)]\label{our:m-cond}
  \uses{our:lsc, coroll:trim-univ, def:m}
Let the condition \( y \) range over \( \dS \), over \( \dZ \), over a space of measures
\( \cM(\bX') \), or over a finite product of these.
 \begin{alphenum}
  \item The \df{conditional universal test} \( \t_{\mu}(x \mvert y) \) is a function, lower
semicomputable in \( (x,y,\mu) \), with \( \mu^{x}\t_{\mu}(x \mvert y)\le 1 \) for all
\( y,\mu \), that dominates every such function up to a multiplicative constant.
It exists by Corollary~\ref{coroll:trim-univ} with \( s = 1 \).
We write \( H_{\mu}(x \mvert y) = -\log\t_{\mu}(x \mvert y) \).
  \item Over \( \dS \), with the counting measure \( \# \) in place of \( \mu \), let
\( \m(x \mvert y) = \t_{\#}(x \mvert y) \): a universal function among those lower
semicomputable in \( (x,y) \) with \( \sum_{x}\m(x \mvert y)\le 1 \) for all \( y \).
Let \( K(x \mvert y) = -\log\m(x \mvert y) = H_{\#}(x \mvert y) \).
In particular \( \m(x \mvert \mu) \) and \( K(x \mvert \mu) \) are defined for every measure
\( \mu \).
 \end{alphenum}
\end{ourdefinition}

\begin{formalization}
Planned. This is how G\'acs uses $K(x \mvert \mu)$ in Section 4.6: as $-\log$ of a semimeasure
that is lower semicomputable in $(x,\mu)$, not as the length of a shortest program (see the
paragraph before Proposition~\ref{propo:H/mu-computable}, and
Remark~\ref{rem:oracle-warning}). In Lean, $\m(x \mvert y)$ is the universal element obtained by
relativizing Chapter 1's trimming (\texttt{exists\_lsc\_semimeasure\_family}, \texttt{truncG})
uniformly in the measure argument. With a string condition it agrees with Chapter 1's conditional
$\m$ up to a constant factor, so by the coding theorem $K(x \mvert y)$ is then the prefix
complexity up to an additive constant. Chapter 3's oracle complexity \texttt{KPrel V A x} for a
name $A$ of $\mu$ satisfies only one inequality (Theorem~\ref{our:oracle-bound}), so it does not
serve as the definition.
\end{formalization}

% Ours: the facts about lower semicomputability that the proofs use (Appendix B)
\begin{ourtheorem}[Lower semicomputability in a measure argument]\label{our:lsc-props}
  \uses{our:lsc}
 \begin{alphenum}
  \item (Oracle form.) A function \( f\ge 0 \) is lower semicomputable if and only if a single
oracle machine, given \( (x,y) \) and any name \( A \) of \( \mu \) as its oracle, enumerates
rationals whose supremum is \( f(x,y,\mu) \).
  \item (Computable approximations, Proposition B.1.38.) A function \( f\ge 0 \) is lower
semicomputable if and only if it is the supremum of a computable sequence of functions
\( f_{i} \), each depending on \( x \) through finitely many strings and on \( \mu \) through
finitely many values \( \mu(u) \), rational-piecewise-linearly.
For such \( f_{i} \), the integral \( \mu^{x}f_{i}(x,y,\mu) \) is computable from \( y \) and a
name of \( \mu \).
  \item (Cantor space, Proposition B.1.23.) Over \( \dB^{\dN} \), \( f \) is lower
semicomputable if and only if \( f(\xi,y,\mu) = \sup_{n}g(\xi^{\le n},y,\mu) \) for a function
\( g \) on strings that is lower semicomputable and monotonic with respect to the prefix
relation.
  \item (Closure.) Lower semicomputability is preserved by suprema of computable sequences, by
finite sums and products of nonnegative functions, by sums \( \sum_{z}\nu(z)f(z,\dots) \) over a
discrete argument against a measure argument \( \nu \), by substituting a computable measure for
a measure argument, by substituting computable maps of measures (the image \( \Lg^{*}\mu \) under
a computable kernel over \( \dS \) whose rows sum to 1, the product \( \mu\times\nu \), the
marginals, the image under a computable map), and by substituting an upper semicomputable
integer-valued function for an integer argument in which the function is decreasing.
 \end{alphenum}
\end{ourtheorem}

\begin{formalization}
Planned, as the toolkit of the chapter. In (a), a prefix of a name confines each value it lists
to a closed rational interval, and every measure in the corresponding open box has a name with
that prefix; the machine reads its oracle monotonically, as Chapter 3's \texttt{IsOracleMachine}
requires. (b) is what makes G\'acs's proof of Theorem~\ref{thm:trim} work: the strict comparison
of a computable integral with a lower semicomputable bound is an open condition on $\mu$. (c) is
how Chapter 2 represents functions on sequences (\texttt{seqVal}, \texttt{IsMonoStr}). In (d) the
image $\Lg^{*}\mu(y) = \sum_{x}\mu(x)\lg_{x}(y)$ is lower semicomputable, and also upper
semicomputable as $1-\sum_{y'\ne y}\Lg^{*}\mu(y')$, the device Chapter 2 uses for image measures
(\texttt{ExactConservation.lean}).
\end{formalization}

% Ours: infima over compact parameter sets (Proposition B.1.31, specialized)
\begin{ourtheorem}[Infima over compact parameter sets]\label{our:lsc-inf}
  \uses{our:lsc}
 \begin{alphenum}
  \item Let \( p\mapsto\mu_{p} \), \( p\in\clint{0}{1} \), be a computable family of measures,
that is \( (w,p)\mapsto\mu_{p}(w) \) is computable, and let \( f \) be lower semicomputable.
Then \( (x,y)\mapsto\inf_{p}f(x,y,\mu_{p}) \) is lower semicomputable.
  \item Let \( u \) be lower semicomputable over \( \dB^{\dN} \), in \( (\xi,\mu) \).
Then \( (z,\mu)\mapsto\inf\setof{u(\og,\mu) : z\prefix\og} \) is lower semicomputable over
\( \dS \).
 \end{alphenum}
\end{ourtheorem}

\begin{formalization}
Planned. These are the two instances of Proposition B.1.31 (the infimum of a lower
semicomputable function over an effectively compact space is lower semicomputable) that the plan
needs: the interval $\clint{0}{1}$, where a finite cover by rational open intervals is
recognizable, and the cylinder of $z$, where a finite cover by cylinders is recognizable. The
Bernoulli family $p\mapsto B_{p}$ of Definition~\ref{def:g4.2.1} is computable, so (a) applies to
it; this replaces Example~\ref{ex:g4.2.3} and Proposition B.1.32 in the plan.
\end{formalization}


% Gacs source line 4994
 \begin{definition}\label{def:mu-test}
  \uses{our:measure, our:lsc, def:g2.2.1, def:g2.3.29}
Let us be given a computable complete metric space \( \bX = (X, d, D, \ag) \).
For an arbitrary measure \( \mu\in\cM(\bX) \), a \( \mu \)-\df{test of randomness}
is a \( \mu \)-lower semicomputable function \( f:X\to\ol\dR_{+} \) with the property
\( \mu f \le 1 \).
We call an element \( x \) \df{random} with respect to \( \mu \)
if \( f(x)<\infty \) for all \( \mu \)-tests \( f \).
But even among random elements, the size of the tests quantifies
(non-)randomness.

A \df{uniform test} of randomness is a lower semicomputable function
\( f:\bX\times\cM(\bX)\to \ol\dR_{+} \), 
written as \( \tup{x,\mu}\mapsto f_{\mu}(x) \)
such that \( f_{\mu}(\cdot) \) is a \( \mu \)-test for each \( \mu \).
 \end{definition}

\begin{formalization}
Planned for $\bX = \dS$ and $\bX = \dB^{\dN}$, with $\cM(\bX)$ and the integral $\mu f$ as in
Definition~\ref{our:measure} and lower semicomputability as in Definition~\ref{our:lsc}. A test
over $\dB^{\dN}$ is represented by a prefix-monotone function on strings
(Theorem~\ref{our:lsc-props}(c)), as Chapter 2 represents its tests (\texttt{IsIntegrableTest},
\texttt{IsIntegrableSeqTest}). For a computable $P$, the $P$-tests are exactly Chapter 2's tests
for $P$.
\end{formalization}

The condition \( \mu f\le 1 \) guarantees that the
probability of those outcomes whose randomness is \( \ge m \) is at most
\( 1/m \).
The definition of tests is in the spirit of Martin-L\"of's tests.
The important difference is in the semicomputability condition:
instead of restricting the measure \( \mu \) to be computable, we require the
test to be lower semicomputable also in its argument \( \mu \).

The following result implies that every \( \mu \)-test can be extended to a uniform
test.

% Gacs source line 5022
 \begin{theorem}\label{thm:trim}
  \uses{our:lsc, our:measure}
Let \( \phi_{e} \), \( e=1,2,\dots \) be an enumeration of all lower semicomputable
functions \( \bX\times\bY\times\cM(\bX)\to\ol\dR_{+} \), where \( \bY \) is also a
computable
metric space, and \( s:\bY\times\cM(\bX)\to\ol\dR \) a lower semicomputable
function.
There is a recursive function \( e\mapsto e' \) with the property that
 \begin{alphenum}
  \item For each \( e \), the function \( \phi_{e'}(x,y,\mu) \) is everywhere defined
with \( \mu^{x}\phi_{e'}(x,y,\mu)\le s(y,\mu) \).     
  \item For each \( e,y,\mu \), if \( \mu^{x} \phi_{e}(x,y\mu)<s(y,\mu) \) then 
\( \phi_{e'}(\cdot,y,\mu)=\phi_{e}(\cdot,y,\mu) \).
 \end{alphenum}
 \end{theorem}
This theorem generalizes a theorem of Hoyrup and Rojas (in allowing a lower
semicomputable upper bound \( s(y,\mu) \)).

\begin{formalization}
Planned, with $\bX$ one of our two spaces and $\bY = \dS$, extended in Section 4.5 to the
conditions of Definition~\ref{our:m-cond}; the counting measure on $\dS$ may replace $\mu^{x}$,
which gives $\m(x \mvert \mu)$. The enumeration $\phi_{e}$ is indexed by codes of computably
enumerable sets of basic sets. G\'acs's proof works as written once Theorem~\ref{our:lsc-props}(b)
supplies the functions $h_{e,i}$, whose integrals $\mu^{x}h_{e,i}(x,y,\mu)$ are computable from a
name of $\mu$. Without a measure argument this is
the trimming formalized for Chapter 1 (\texttt{truncG} in
\texttt{exists\_lsc\_semimeasure\_family}) and, under the constraints of a continuous
semimeasure, for Theorem~2.3.5 in Chapter 2 (\texttt{ContTrim.lean}).
\end{formalization}
 \begin{proof}
  \proves{thm:trim}
  \uses{our:lsc-props, thm:univsem, thm:g2.3.5}
By Proposition~B.1.38, we can represent \( \phi_{e}(x,y,\mu) \) 
as a supremum 
\( \phi_{e}=\sup_{i}h_{e,i} \) where \( h_{e,i}(x,y,\mu) \) is a computable function of
\( e,i \) monotonically increasing in \( i \).
Similarly, we can represent \( s(y,\mu) \) as a supremum \( \sup_{j}s_{j}(y,\mu) \)
where \( s_{j}(y,\mu) \) is a computable function monotonically increasing in \( j \).
The integral \( \mu^{x} h_{e,i}(x,y,\mu) \) is computable as a function of \( (y,\mu) \),
in particular it is upper semicomputable.

Define \( h'_{e,i,j}(x,y,\mu) \) as \( h_{e,i}(x,y,\mu) \) for all \( j,y,\mu \)
with \( \mu^{x} h_{e,i}(x,y,\mu)<s_{j}(y,\mu) \), and 0 otherwise.
Since \( s_{j}(y,\mu) \) is computable this 
definition makes the function \( h'_{e,i,j}(x,y,\mu) \) lower semicomputable.
The function \( h''_{e,i}(x,y,\mu)=\sup_{j}h'_{e,i,j}(x,y,\mu) \) is then also lower
semicomputable, with \( h''_{e,i}(x,y,\mu)\le h_{e,i}(x,y,\mu) \),
and \( \mu^{x} h'_{e,i}(x,y,\mu)\le s(y,\mu) \).
Also, \( h''_{e,i}(x,y,\mu) \) is monotonically increasing in \( i \).
The function \( \phi'_{e}(x,y,\mu)=\sup_{i}h''_{e,i}(x,y,\mu) \) is then also lower
semicomputable, and by Fubini's theorem we have
\( \mu^{x}\phi'_{e}(x,y,\mu)\le s(y,\mu) \).

Define \( \phi_{e'}(x,y,\mu)=\phi'_{e}(x,y,\mu) \).
Consider any \( e,y,\mu \) such that \( \mu^{x}\phi_{e}(x,y,\mu)<s(y,\mu) \) holds.
Then for every \( i \) there is a \( j \) with \( \mu^{x}h_{e,i}(x,y,\mu)<s_{j}(y,\mu) \),
and hence \( h'_{e,i,j}(x,y,\mu)=h_{e,i}(x,y,\mu) \).
It follows that \( h''_{e,i}(x,y,\mu)=h_{e,i}(x,y,\mu) \) for all \( i \) and hence
\( \phi'_{e}(x,y,\mu)=\phi_{e}(x,y,\mu) \).
 \end{proof}

% Gacs source line 5074
\begin{corollary}[Uniform extension]\label{coroll:trim}
  \uses{def:mu-test}
There is an operation \( H_{e}(x,\mu)\mapsto H_{e'}(x,\mu) \) with the property
that \( H_{e'}(x,\mu) \) is a uniform test and if \( 2 H_{e}(\cdot,\mu) \) is a \( \mu \)-test
then \( H_{e'}(x,\mu)=H_{e}(x,\mu) \).
\end{corollary}
 \begin{proof}
  \proves{coroll:trim}
  \uses{thm:trim}
In Theorem~\ref{thm:trim} set \( \phi_{e}(x,y,\mu)=\frac{1}{2}H_{e}(x,\mu) \) with
\( s(y,\mu)=1 \).
 \end{proof}

% Gacs source line 5084
\begin{corollary}[Universal generalized test]\label{coroll:trim-univ}
  \uses{our:lsc}
Let \( s:\bY\times\cM(\bX)\to\ol\dR_{+} \) a lower semicomputable
function.
Let \( E \) be the set of 
lower semicomputable functions \( \phi(x,y,\mu)\ge 0 \) with
\( \mu^{x}\phi(x,y,\mu)\le s(y,\mu) \).
There is a function \( \psi\in E \) that is \emph{optimal} in the sense
that for all \( \phi\in E \) there is a constant \( c_{\phi} \) with 
\( \phi\le 2c_{\phi}\psi \).
\end{corollary}
 \begin{proof}
  \proves{coroll:trim-univ}
  \uses{thm:trim}
Apply the operation \( e\mapsto e' \) of theorem~\ref{thm:trim} to the sequence
\( \phi_{e}(x,y,\mu) \) (\( e=1,2,\dots \)) of all lower semicomputable
functions of \( x,y,\mu \).
The elements of the sequence \( \phi'_{e}(x,y,\mu) \), \( e=1,2,\dots \)  are in \( E \) and
and the sequence \( 2\phi'_{e}(x,y,\mu) \), \( e=1,2,\dots \) contains all
elements of \( E \).
Hence 
the function \( \psi(x,y,\mu)=\sum_{e=1}^{\infty}2^{-e} \phi'_{e}(x,y,\mu) \) is in
\( E \) and has the optimality property.
 \end{proof}

% Gacs source line 5106
 \begin{definition}
  \label{def:g4.1.4}
  \uses{def:mu-test}
A uniform test \( u \) is called \df{universal} if for every other test \( t \) there
is a constant \( c_{t}>0 \) such that for all \( x,\mu \) we have \( t_{\mu}(x) \le c u_{\mu}(x) \). 
 \end{definition}

% Gacs source line 5111
\begin{theorem}[Universal test,\cite{HoyrupRojasRandomness09}]
  \label{thm:g4.1.2}
  \uses{def:g4.1.4, thm:univ-integrable-test, thm:univ-integrable-test-seq}
There is a universal uniform test.
\end{theorem}

\begin{formalization}
Planned. Over $\dS$ the universal test can be taken to be $\m(x \mvert \mu)/\mu(x)$, and its
universality is Theorem~\ref{thm:test-charac-discr}, as Chapter 2 did for a computable $P$
(\texttt{canonicalIntegrableTest\_isUniversal}); only the existence of $\m(x \mvert \mu)$ needs
trimming. Over $\dB^{\dN}$ the trimming of Theorem~\ref{thm:trim} is needed: Chapter 2's explicit
universal test (\texttt{canonicalSeqTest\_isUniversal}) has no counterpart for noncomputable
$\mu$, since Proposition~\ref{propo:test-charac-seq-lb} gives only one inequality.
\end{formalization}
 \begin{proof}
  \proves{thm:g4.1.2}
  \uses{coroll:trim-univ}
This is a special case of Corollary~\ref{coroll:trim-univ} with \( s(y,\mu)=1 \).
 \end{proof}

% Gacs source line 5118
\begin{definition}
  \label{def:g4.1.5}
  \uses{thm:g4.1.2}
Let us fix a universal uniform test, called \( \t_{\mu}(x) \).
An element \( x\in X \) is called \df{random} with respect to measure
\( \mu\in\cM(\bX) \) if \( \t_{\mu}(x)<\infty \).

The \df{deficiency of randomness} is defined as \( \d_{\mu}(x) = \log\t_{\mu}(x) \).
\end{definition}

If the space is discrete then typically all elements are random with respect to
\( \mu \), but they will still be distinguished according to their different values
of \( \d_{\mu}(x) \).

\subsection{Sequences}

Let our computable metric space \( \bX=(X,d,D,\ag) \) be the Cantor space
of Example~B.1.35.2:
the set of sequences over a (finite or countable) alphabet \( \Sg^{\dN} \).
We may want to measure the non-randomness of finite sequences, viewing them
as initial segments of infinite sequences.
Take the universal test \( \t_{\mu}(x) \).
For this, it is helpful to apply the representation of
Proposition~B.1.23, taking into account that adding
the extra parameter \( \mu \) does not change the validity of the theorem:

% Gacs source line 5142
 \begin{proposition}\label{propo:finite-test}
  \uses{def:g4.1.5, our:lsc-props}
There is a function
\( g:\cM(\bX)\times\Sg^{*}\to\ol\dR_{+} \) 
with \( \t_{\mu}(\xi) = \sup_{n} g_{\mu}(\xi^{\le n}) \), and 
with the following properties:
 \begin{alphenum}
  \item \( g \) is lower semicomputable.
  \item \( v\prefix w \) implies \( g_{\mu}(v)\le g_{\mu}(w) \).
  \item\label{i:extended.sum} For all integer \( n\ge 0 \) we have 
\( \sum_{w\in\Sg^{n}}\mu(w)g_{\mu}(w) \le 1 \).
 \end{alphenum}
 \end{proposition}

\begin{formalization}
Planned. Tests over $\dB^{\dN}$ are given by prefix-monotone functions on strings from the start
(Definition~\ref{def:mu-test}), so this is Theorem~\ref{our:lsc-props}(c) for $\t$, with (c) the
level form of $\mu\t_{\mu}\le 1$, as in Chapter 2's \texttt{IsIntegrableSeqTest}.
\end{formalization}
The properties of the function \( g_{\mu}(w) \) clearly imply that 
\( \sup_{n} g_{\mu}(\xi^{\le n}) \) is a uniform test.

The existence of a universal function among the functions \( g \) can be proved by
the usual methods:

% Gacs source line 5160
 \begin{proposition}\label{propo:finite-univ}
  \uses{propo:finite-test, thm:trim, thm:g2.3.5}
Among the functions \( g_{\mu}(w) \) satisfying the properties listed in
Proposition~\ref{propo:finite-test}, there is one that
dominates to within a multiplicative constant.
 \end{proposition}

\begin{formalization}
Planned: trimming as in Theorem~\ref{thm:trim}, now to the constraints (c) at every level $n$,
as Chapter 2 trims to the constraints of a continuous semimeasure (\texttt{ContTrim.lean}).
\end{formalization}

These facts motivate the following definition.

% Gacs source line 5168
 \begin{definition}[Extended test]\label{def:extended-test}
  \uses{propo:finite-univ, def:g2.3.29}
Over the Cantor space, we
extend the definition of a universal test \( \t_{\mu}(x) \) to finite sequences
as follows.
We fix a function \( \t_{\mu}(w) \) with \( w\in\Sg^{*} \)
whose existence is assured by~Proposition~\ref{propo:finite-univ}.
For infinite sequences \( \xi \) we define 
\( \t_{\mu}(\xi) = \sup_{n}\t_{\mu}(\xi^{\le n}) \).
The test with values defined also on finite sequences will be called an
\df{extended test}.
 \end{definition}

We could try to define extended tests also over arbitrary constructive
metric spaces, extending them
to the canonical balls, with the monotonicity property that
\( \t_{\mu}(v)\le\t_{\mu}(w) \) if ball \( w \) is manifestly included in ball \( v \).
But there is nothing simple and obvious corresponding to the 
integral requirement~\eqref{i:extended.sum}.

Over the space \( \Sg^{\dN} \) for a finite alphabet \( \Sg \),
an extended test could also be extracted directly from a test, 
using the following formula (as observed by Vyugin and Shen).

% Gacs source line 5191
\begin{definition}\label{def:ol-u}
  \uses{our:lsc-inf}
 \begin{align*}
   \ol\u(z) = \inf\setof{\u(\og) : \og\text{ is an infinite extension of } z}.
 \end{align*}
\end{definition}
This function is lower semicomputable, by Proposition~B.1.31.

\begin{formalization}
Planned for $\Sg = \dB$ and uniformly in $\mu$; the lower semicomputability is
Theorem~\ref{our:lsc-inf}(b).
\end{formalization}


\subsection{Conservation of randomness}

For \( i=1,0 \), let
\( \bX_{i} = (X_{i}, d_{i}, D_{i}, \ag_{i}) \) be computable metric spaces, 
and let \( \bM_{i} = (\cM(\bX_{i}), \sg_{i}, \nu_{i}) \) 
be the effective topological space of
probability measures over \( \bX_{i} \).
Let \( \Lg \) be a computable probability kernel from \( \bX_{1} \) to \( \bX_{0} \) as
defined in Subsection~B.2.3.
In the following theorem, the same notation \( \d_{\mu}(x) \) will
refer to the deficiency of randomness with respect to two different spaces,
\( \bX_{1} \) and \( \bX_{0} \), but this should not cause confusion.
Let us first spell out the conservation theorem before interpreting it.

% Gacs source line 5218
 \begin{theorem}\label{propo:conservation}
  \uses{def:g4.1.5, thm:g2.2.3}
For a computable probability kernel \( \Lg \) from \( \bX_{1} \) to \( \bX_{0} \),
we have 
 \begin{equation}\label{eq:conservation}
  \lg_{x}^{y} \t_{\Lg^{*}\mu}(y) \lem \t_{\mu}(x).
 \end{equation}
 \end{theorem}

\begin{formalization}
Planned for computable kernels from $\dS$ to $\dS$: Chapter 2's \texttt{ComputableTransition}
($\lg_{x}(y)$ lower semicomputable, rows summing to 1), now with an arbitrary input measure
$\mu$. This generalizes Theorem~2.2.3 (\texttt{transTest\_le\_canonical}), where $\mu$ is
computable. That $u_{\mu}(x)$ is lower semicomputable is Theorem~\ref{our:lsc-props}(d). Kernels
into or out of $\dB^{\dN}$ are not planned, except the deterministic ones of
Corollary~\ref{coroll:conservation-determ}.
\end{formalization}
 \begin{proof}
  \proves{propo:conservation}
  \uses{thm:g4.1.2, our:lsc-props}
Let \( \t_{\nu}(x) \) be the universal test over \( \bX_{0} \).
The left-hand side of~\eqref{eq:conservation} can be written as
 \[
  u_{\mu} = \Lg \t_{\Lg^{*}\mu}.
 \]
According to~(A.2.4), we have 
\( \mu u_{\mu} = (\Lg^{*}\mu) \t_{\Lg^{*}\mu} \) which is \( \le 1 \) since \( \t \) is
a test.
If we show that \( (\mu,x) \mapsto u_{\mu}(x) \) is lower semicomputable
then the universality of \( \t_{\mu} \) will imply
that \( u_{\mu} \lem \t_{\mu} \).

According to Proposition~B.1.38, as
a lower semicomputable function, \( \t_{\nu}(y) \) can be written as 
\( \sup_{n} g_{n}(\nu, y) \), where \( (g_{n}(\nu, y)) \) is a computable sequence
of computable functions.
We pointed out in Subsection~B.2.3 that the function 
\( \mu \mapsto \Lg^{*}\mu \) is computable.
Therefore the function \( (n, \mu, x) \mapsto g_{n}(\Lg^{*}\mu, f(x)) \) 
is also a computable.
So, \( u_{\mu}(x) \) is the supremum of a computable sequence of computable
functions and as such, lower semicomputable.
 \end{proof}

It is easier to interpret the theorem first in the special case when
\( \Lg = \Lg_{h} \) for a computable function \( h : X_{1} \to X_{0} \),
as in Example~B.2.12.
Then the theorem simplifies to the following.

% Gacs source line 5256
 \begin{corollary}\label{coroll:conservation-determ}
  \uses{propo:conservation, thm:simple-rand-cons-det}
For a computable function \( h : X_{1} \to X_{0} \), we have
\( \d_{h^{*}\mu}(h(x)) \lea \d_{\mu}(x) \).
 \end{corollary}

\begin{formalization}
Planned for computable $h:\dS\to\dS$ (the kernel of Chapter 2's \texttt{detTransition}) and for
computable $h:\dB^{\dN}\to\dB^{\dN}$, given by a monotone map on strings. In the second case
$h^{*}\mu(w)$ is a finite sum of values of $\mu$, by compactness, and the proof of
Theorem~\ref{propo:conservation} applies with level integrals. This generalizes
Theorem~2.2.2 (\texttt{canonical\_pushforward\_le}, \texttt{det\_canonical\_image\_le}).
\end{formalization}

Informally, this says that if \( x \) is random with respect to \( \mu \) in
\( \bX_{1} \) then \( h(x) \) is essentially at least as random with respect to the
output distribution \( h^{*}\mu \) in \( \bX_{0} \).
Decrease in randomness
can only be caused by complexity in the definition of the function \( h \).

Let us specialize the theorem even more:

% Gacs source line 5270
\begin{corollary}\label{coroll:conservation-pair}
  \uses{coroll:conservation-determ}
For a probability distribution \( \mu \) over the space \( X\times Y \)
let \( \mu_{X} \) be its marginal on the space \( X \).
Then we have
 \begin{align*}
   \d_{\mu_{X}}(x) \lea \d_{\mu}(\tup{x,y}).
 \end{align*}
\end{corollary}

\begin{formalization}
Planned over $\dS$, with pairs coded by \texttt{pairCode}, and over $\dB^{\dN}$, with pairs of
sequences interleaved. In both cases the projection is a computable map, and this is
Corollary~\ref{coroll:conservation-determ} for it.
\end{formalization}

This says, informally, that if a pair is random then each of its elements is
random (with respect to the corresponding marginal distribution).

In the general case of the theorem, concerning random transitions,
we cannot bound the randomness of each outcome uniformly.
The theorem asserts that the average nonrandomness, as measured by 
the universal test with respect to the output distribution, does not
increase.
In logarithmic notation:
\( \lg_{x}^{y} 2^{\d_{\Lg^{*}\mu}(y)} \lea \d_{\mu}(x) \), 
or equivalently, 
\( \int 2^{\d_{\Lg^{*}\mu}(y)} \lg_{x}(dy) \lea \d_{\mu}(x) \).

% Gacs source line 5292
 \begin{corollary}
  \label{cor:g4.1.12}
  \uses{propo:conservation, cor:g2.2.8}
Let \( \Lg \) be a computable probability kernel from \( \bX_{1} \) to \( \bX_{0} \).
There is a constant \( c \) such that
for every \( x\in\bX^{1} \), and integer \( m > 0 \) we have 
 \[
 \lg_{x}\setof{y : \d_{\Lg^{*}\mu}(y) > \d_{\mu}(x) + m + c} \le 2^{-m}.
 \]
 \end{corollary}

\begin{formalization}
Planned for the kernels of Theorem~\ref{propo:conservation}. As in Chapter 2's
Corollary~2.2.8 (\texttt{transition\_deficiency\_tail}), it follows by Markov's inequality.
\end{formalization}

Thus, in a computable random transition,
the probability of an increase of randomness deficiency by
\( m \) units (plus a constant \( c \)) is less than \( 2^{-m} \).
The constant \( c \) comes from the description complexity 
of the transition \( \Lg \).

A randomness conservation result related to 
Corollary~\ref{coroll:conservation-determ} was proved 
in~\cite{HertlingWeihrauchRand98}.
There, the measure over the space \( \bX_{0} \) is not the output
measure of the transformation, but is assumed to
obey certain inequalities related to the transformation.


\section{Test for a class of measures}

\subsection{From a uniform test}

A Bernoulli measure is what we get by tossing a (possibly biased) coin
repeatedly.

% Gacs source line 5322
\begin{definition}
  \label{def:g4.2.1}
  \uses{our:measure}
Let \( X=\dB^{\dN} \) be the set of infinite binary sequences, with the usual
sequence topology.
Let \( B_{p} \) be the measure on \( X \) that corresponds to tossing a coin
independently with probability \( p \) of success: it is called the \df{Bernoulli}
measure with parameter \( p \).
Let \( \cB \) denote the set of all Bernoulli measures. 
\end{definition}

\begin{formalization}
Planned: $B_{p}(w) = p^{k}(1-p)^{n-k}$ for a string $w$ of length $n$ with $k$ ones, a measure over
$\dB^{\dN}$ in the sense of Definition~\ref{our:measure}, computable in $(w,p)$.
\end{formalization}

Given a sequence \( x\in X \) we may ask the question whether \( x \) is random with
respect to at least \emph{some} Bernoulli measure.
(It can clearly not be random with respect to two different Bernoulli measures
since if \( x \) is random with respect to \( B_{p} \) then its relative frequencies
converge to \( p \).)
This idea suggests two possible definitions for a test of the property of
``Bernoulliness'': 
 \begin{enumerate}
  \item We could define \( \t_{\cB}(x) = \inf_{\mu\in\cB}(x) \).
  \item We could define the notion of a Bernoulli test as a lower semicomputable
    function \( f(x) \) with the property \( B_{p} f\le 1 \) for all \( p \).
 \end{enumerate}
We will see that in case of this class of measures the two definitions lead to
essentially the same test.

Let us first extend the definition to more general sets of measures, still
having a convenient property.
 
% Gacs source line 5349
 \begin{definition}[Class tests]\label{def:C-test}
  \uses{def:mu-test, our:lsc}
Consider a class \( \cC \) of measures that is effectively
compact in the sense of Definition~B.1.27 or (equivalently
for metric spaces) in the sense of Theorem~B.1.1.
A lower semicomputable function \( f(x) \) is called a \( \cC \)-\df{test} if for all
\( \mu\in\cC \) we have \( \mu f\le 1 \).
It is a \df{universal} \( \cC \)-\df{test} if it dominates all other \( \cC \)-tests to
within a multiplicative constant.
\end{definition}

\begin{formalization}
Planned for classes that are computable images $\cC = \setof{\mu_{p} : p\in\clint{0}{1}}$ of the
unit interval, that is with $(w,p)\mapsto\mu_{p}(w)$ computable. These include $\cB$ and the
Markov chain classes of Example~\ref{ex:g4.2.9}, and they are effectively compact by
Proposition B.1.32; the plan uses no other kind of class. General effectively compact classes of
measures are out of scope.
\end{formalization}

% Gacs source line 5359
\begin{example}
  \label{ex:g4.2.3}
It is easy to show that the class \( \cB \) is effectively compact.
One way is to appeal to the general theorem in
Proposition~B.1.32 saying that applying a computable
function to an effectively compact set (in this case the interval
\( \clint{0}{1} \)), the image is also an effectively compact set.
\end{example}

For the case of infinite sequences, we can also define extended tests.

% Gacs source line 5369
 \begin{definition}[Extended class test]
  \label{def:g4.2.4}
  \uses{def:C-test, def:extended-test}
Let our space \( \bX \) be the Cantor space of infinite sequences \( \Sg^{\dN} \).
Consider a class \( \cC \) of measures that is effectively
compact in the sense of Definition~B.1.27 or (equivalently
for metric spaces) in the sense of Theorem~B.1.1.
A lower semicomputable function \( f:\Sg^{*}\to\ol\dR_{+} \) 
is called an \df{extended} \( \cC \)-\df{test} if it is monotonic with respect to
the prefix relation and for all \( \mu\in\cC \) and integer \( n\ge 0 \) we have 
 \begin{align*}
  \sum_{x\in\Sg^{n}}\mu(x) f(x) \le 1.   
 \end{align*}
It is \df{universal} if it dominates all other extended \( \cC \)-tests to
within a multiplicative constant.
\end{definition}

The following observation is immediate.

% Gacs source line 5386
 \begin{proposition}
  \label{prop:g4.2.5}
  \uses{def:g4.2.4, our:lsc-props}
A function \( f:\Sg^{\dN}\to\ol\dR_{+} \) is a class test if and only if it can be
represented as \( \lim_{n}g(\xi^{\le n}) \) where \( g(x) \) is an extended class test.
 \end{proposition}

\begin{formalization}
Planned: Theorem~\ref{our:lsc-props}(c) and monotone convergence, as in Chapter 2
(\texttt{seqVal}).
\end{formalization}

The following theorem defines a universal \( \cC \)-test.

% Gacs source line 5394
 \begin{theorem}
  \label{thm:g4.2.1}
  \uses{def:C-test, def:g4.1.5}
Let \( \t_{\mu}(x) \) be a universal uniform test.
Then \( \u(x) = \inf_{\mu\in\cC}\t_{\mu}(x) \) defines a universal \( \cC \)-test.
 \end{theorem}

\begin{formalization}
Planned for the classes of Definition~\ref{def:C-test}; the infimum is lower semicomputable by
Theorem~\ref{our:lsc-inf}(a).
\end{formalization}
 \begin{proof}
  \proves{thm:g4.2.1}
  \uses{coroll:trim, our:lsc-inf}
Let us show first that \( \u(x) \) is a \( \cC \)-test.
It is lower semicomputable according to Proposition~B.1.31.
Also, for each \( \mu \) we have \( \mu\u \le \mu \t_{\mu} \le 1 \), showing that
\( \u(x) \) is a \( \cC \)-test.

Let us now show that it is universal.
Let \( f(x) \) be an arbitrary \( \cC \)-test.
By Corollary~\ref{coroll:trim} there is a uniform test \( g_{\mu}(x) \) such that for
all \( \mu\in\cC \) we have \( g_{\mu}(x)=f(x)/2 \).
It follows from the universality of the uniform test \( \t_{\mu}(x) \)
that \( f(x) \lem g_{\mu}(x) \lem\t_{\mu}(x) \) for all \( \mu\in\cC \).
But then \( f(x) \lem\inf_{\mu\in\cC}\t_{\mu}(x) = \u(x) \).
 \end{proof}

For the case of sequences, the same statement can be made for extended tests.
(This is not completely automatic since a test is obtained from an
extended test via a supremum, on the other hand a class test is obtained,
according the theorem above, via an infimum.)


\subsection{Typicality and class tests}

The set of Bernoulli measures has an important property shared by many
classes considered in practice: namely that random sequences determine the
measure to which it belongs.

% Gacs source line 5426
 \begin{definition}
  \label{def:g4.2.6}
  \uses{def:C-test}
Consider a class \( \cC \) of measures over a computable metric space
\( \bX = (X,d,D,\ag) \).
We will say that a lower semicomputable function
 \begin{align*}
  \s:X\times\cC\to\ol\dR_{+}
 \end{align*}
is a \df{separating test} for \( \cC \) if 
\begin{itemize}
\item \( \s_{\mu}(\cdot) \) is a test for each \( \mu\in\cC \).
\item If \( \mu\ne\nu \) then
\( \s_{\mu}(x)\lor\s_{\nu}(x)=\infty \) for all \( x\in X \).
\end{itemize}
Given a separating test \( \s_{\mu}(x) \) we call an element \( x \)
\df{typical} for \( \mu\in\cC \) if \( \s_{\mu}(x)<\infty \).
 \end{definition}

A typical element determines uniquely the measure \( \mu \) for which
it is typical.
Note that if a separating tests exists for a class then any two different
measures \( \mu_{1},\mu_{2} \) in the class are orthogonal to each other, that is there are
disjoint measureable sets \( A_{1},A_{2} \) with \( \mu_{j}(A_{j})=1 \).
Indeed, let \( A_{j}=\setof{x:s_{\mu_{j}}(x)<\infty} \).

Let us show a nontrivial example: the class of \( \cB \) of Bernoulli measures.
Recall that by Chebyshev's inequality we have
 \begin{align*}
   B_{p}(\setof{x\in\dB^{n}: |\sum_{i}x(i)-n p|\ge\lg n^{1/2}(p(1-p))^{1/2}}) \le \lg^{-2}.
 \end{align*}
Since \( p(1-p)\le 1/4 \), this implies
 \begin{align*}
   B_{p}(\setof{x\in\dB^{n}: |\sum_{i}x(i)-n p>\lg n^{1/2}/2}) < \lg^{-2}.
 \end{align*}
Setting \( \lg=n^{0.1} \) and ignoring the factor \( 1/2 \) gives
 \begin{align*}
   B_{p}(\setof{x\in\dB^{n}: |\sum_{i}x(i)-n p|> n^{0.6}})< n^{-0.2}.
 \end{align*}
Setting \( n=2^{k} \):
 \begin{align}\label{eq:2^k-Cheb}
   B_{p}(\setof{x\in\dB^{2^{k}}: |\sum_{i}x(i)-2^{k}p|> 2^{0.6 k}})< 2^{-0.2k}.
 \end{align}

Now, for the example.

\begin{example}\label{example:Bernoulli-sep}
For a sequence \( \xi \) in \( \bB^{\dN} \), and for \( p\in\clint{0}{1} \) let
 \begin{align*}
   g_{p}(x) = g_{B_{p}}(x) = \sup\setof{k: |\sum_{i=1}^{2^{k}}\xi(i)-2^{k}p|> 2^{0.6k}}.
 \end{align*}
Then we have
 \begin{align*}
   B_{p}^{\xi}g_{p}(\xi) \le \sum_{k} k\cdot 2^{-0.2 k} = c<\infty
 \end{align*}
for some constant \( c \), so \( \s_{p}(\xi)=g_{p}(x)/c \) is a test for each \( p \).
The property \( \s_{p}(\xi)<\infty \) implies that \( 2^{-k}\sum_{i=1}^{2^{k}}\xi(i) \)
converges to \( p \).
For a given \( \xi \) this is impossible for both \( p \) and \( q \) for \( p\ne q \), hence
\( \s_{p}(\xi)\lor\s_{q}(\xi)=\infty \).
\end{example}

The following structure theorem gives considerable insight.

% Gacs source line 5488
\begin{theorem}\label{thm:test-sep}
  \uses{def:g4.2.6, thm:g4.2.1}
Let \( \cC \) be an effectively compact class of measures,
let \( \t_{\mu}(x) \) be the universal uniform test and
let \( \t_{\cC}(x) \) be a universal class test for \( \cC \).
Assume that a separating test \( \s_{\mu}(x) \) exists for \( \cC \).
Then we have the representation
 \begin{align*}
  \t_{\mu}(x) \eqm \t_{\cC}(x)\lor \s_{\mu}(x)
 \end{align*}
for all \( \mu\in\cC \), \( x\in X \).
\end{theorem}

\begin{formalization}
Planned, with the Bernoulli instance: the separating test of Example~\ref{example:Bernoulli-sep}
rests on Chebyshev's inequality~\eqref{eq:2^k-Cheb} and is lower semicomputable in $(\xi,p)$, since its
condition on $p$ is a strict inequality.
\end{formalization}
 \begin{proof}
  \proves{thm:test-sep}
  \uses{coroll:trim, thm:g4.2.1}
First, we have \( \t_{\cC}(x)\lor\s_{\mu}(x) \lem \t_{\mu}(x) \).
Indeed as we know from the Uniform Extension Corollary~\ref{coroll:trim}, we can
extend \( \s_{\mu}(x)/2 \) to a uniform test, hence \( \s_{\mu}(x)\lem \t_{\mu}(x) \).
Also by definition \( \t_{\cC}(x)\le \t_{\mu}(x) \).

On the other hand, let us show \( \t_{\cC}(x)\lor\s_{\mu}(x) \ge \t_{\mu}(x) \).
Suppose first that 
\( x \) is not random with respect to any \( \nu\in\cC \): then \( \t_{\cC}(x)=\infty \).
Suppose now that \( x \) is random with respect to some \( \nu\in\cC \), 
\( \nu\ne\mu \).
Then \( \s_{\mu}(x)=\infty \).
Finally, suppose \( \t_{\mu}(x)<\infty \).
Then \( \t_{\nu}(x)=\infty \) for all \( \nu\in\cC \), \( \nu\ne\mu \), hence
\( \t_{\cC}(x) = \inf_{\nu\in\cC}\t_{\nu}(x) = \t_{\mu}(x) \), so
the inequality holds again.
 \end{proof}

The above theorem separates the randomness test into two parts.
One part tests randomness with respect to the class \( \cC \), the other one tests
typicality with respect to the measure \( \mu \).
In the Bernoulli example,
 \begin{itemize}
  \item Part \( \t_{\cB}(\xi) \) checks ``Bernoulliness'', that is independence.
It encompasses all the irregularity criteria.
  \item Part \( \s_{p}(\xi) \) checks
(crudely) for the law of large numbers: whether relative frequency converges
(fast) to \( p \).
 \end{itemize}
If the independence of the sequence is taken for granted, we may assume that
the class test is satisfied.
What remains is typicality testing, which is similar to 
ordinary statistical parameter testing.

% Gacs source line 5533
 \begin{remarks}
  \label{rem:g4.2.8}
   \begin{enumerate}
   \item 
Separation is the only requirement of the test \( \s_{\mu}(x) \),
otherwise, for example in the Bernoulli test case,
no matter how crude the convergence criterion expressed by
\( \s_{\mu}(x) \), the maximum \( \t_{\cC}(x)\lor\s_{\mu}(x) \) 
is always (essentially) the same universal test.
  \item
Though the convergence criterion can be crude, but one still seems to need some
kind of constructive convergence of the relative frequencies if the separation
test is to be defined in terms of relative frequency convergence.
   \end{enumerate}
 \end{remarks}

% Gacs source line 5548
 \begin{example}
  \label{ex:g4.2.9}
For \( \eps>0 \) let \( P(\eps) \) be the Markov chain \( X_{1},X_{2},\dots \)
with set of states \( \{0,1\} \),
with transition probabilities \( T(0,1)=T(1,0)=\eps \) and \( T(0,0)=T(1,1)=1-\eps \),
and with \( P\evof{X_{1}=0}=P\evof{X_{1}=1}=1/2 \).
Let \( \cC(\dg) \) be the class of all \( P(\eps) \) with \( \eps\ge\dg \).
For each \( \dg>0 \), this is an effectively compact class, and a separating test is easy to
construct since an effective law of large numbers holds for these Markov
chains. 

We can generalize to the set of \( m \)-state stationary Markov chains whose 
eigenvalue gap is \( \ge\eps \).
 \end{example}

This example is in contrast to V'yugin's example~\cite{VyuginErgodic98}
showing that, in the nonergodic case,
in general no recursive speed of convergence can be guaranteed in the
Ergodic Theorem (which is the appropriate generalization of the law of large
numbers)\footnote{This paper of Vyugin also shows that though there is no recursive speed of
convergence, a certain constructive 
proof of the pointwise ergodic theorem still gives rise to a test of
randomness.}.

We can show that if a computable Markov chain is ergodic then
its law of large numbers \emph{does} have a constructive speed of convergence.
Hopefully, this observation can be extended to some interesting 
compact classes of ergodic processes.


\subsection{Martin-L\"of's approach}

Martin-L\"of also gave a definition of Bernoulli tests in~\cite{MLof66art}.
For its definition let us introduce the following notation.

 \begin{notation} The set of sequences with a given frequency of 1's will be
   denoted as follows:
 \begin{align*}
   \dB(n,k) = \setof{x\in\dB^{n}: \sum_{i}x(i)=k}.
 \end{align*}
 \end{notation}

Martin-L\"of's definition translates to the integral constraint version as follows:

% Gacs source line 5591
 \begin{definition}\label{def:combinat-Bernoulli}
  \uses{def:semicomputability}
Let \( X=\dB^{\dN} \) be the set of infinite binary sequences with the usual
metrics.
A \df{combinatorial Bernoulli test} is a function
\( f:\dB^{*}\to\ol\dR_{+} \) with the following constraints:
 \begin{alphenum}
  \item It is lower semicomputable.
  \item It is monotonic with respect to the prefix relation.
  \item For all \( 0\le k\le n \) we have
 \begin{align}\label{eq:combinat-Bernoulli-sum}
      \sum_{x\in\dB(n,k)} f(x) \le \binom{n}{k}.
 \end{align}
 \end{alphenum}
 \end{definition}

\begin{formalization}
Planned, on strings only: Chapter 2's unary lower semicomputability and \texttt{IsMonoStr}; the
sum over $\dB(n,k)$ is finite.
\end{formalization}

The following observation is useful.

% Gacs source line 5608
 \begin{proposition}\label{propo:Bernoulli-extend}
  \uses{def:combinat-Bernoulli}
If a combinatorial Bernoulli test \( f(x) \) is given on strings \( x \)
of length less than \( n \), then extending it to longer strings using 
monotonicity we get a function that is still a combinatorial Bernoulli test.
 \end{proposition}
 \begin{proof}
  \proves{propo:Bernoulli-extend}
It is sufficient to check the relation~\eqref{eq:combinat-Bernoulli-sum}.
We have (even for \( k=0 \), when \( \dB(n-1,k-1)=\emptyset \)):
 \begin{align*}
  \sum_{x\in\dB(n,k)} f(x) &\le
  \sum_{y\in\dB(n-1,k-1)}f(y)+\sum_{y\in\dB(n-1,k)} f(y) 
\\ &\le \binom{n-1}{k-1} + \binom{n-1}{k} = \binom{n}{k}.
 \end{align*}
 \end{proof}

The following can be shown using standard methods:
  
% Gacs source line 5625
 \begin{proposition}[Universal combinatorial Bernoulli test]
  \label{prop:g4.2.13}
  \uses{def:combinat-Bernoulli, propo:Bernoulli-extend, thm:univsem}
There is a universal combinatorial Bernoulli test \( f(x) \), that is a combinatorial
Bernoulli test with the property that for every
combinatorial Bernoulli test \( h(x) \) there is a constant \( c_{h}>0 \) such that for
all \( x \) we have \( h(x) \le c_{h} g(x) \).
 \end{proposition}

\begin{formalization}
Planned, by the standard method: trim every enumerated lower semicomputable function stage by
stage to the constraints~\eqref{eq:combinat-Bernoulli-sum} on strings of the current length, extend it to longer strings by
monotonicity (Proposition~\ref{propo:Bernoulli-extend}), and mix with weights $2^{-e}$, as
Chapter 1 trims to $\sum\le 1$.
\end{formalization}

% Gacs source line 5632
\begin{definition}
  \label{def:g4.2.14}
  \uses{prop:g4.2.13, thm:g4.2.1}
Let us fix a universal combinatorial Bernoulli test \( b(x) \) and extend it to
infinite sequences \( \xi \) by
 \begin{align*}
   b(\xi) = \sup_{n} b(\xi^{\le n}).
 \end{align*}
Let \( \t_{\cB}(\xi) \) be a universal class test for Bernoulli measures, for
\( \xi\in\dB^{\dN} \).
\end{definition}

Let us show that the general class test for Bernoulli measures and
Martin-L\"of's Bernoulli test yield the same random sequences.

% Gacs source line 5645
 \begin{theorem}\label{thm:combinat-Bernoulli}
  \uses{def:g4.2.14}
With the above definitions, we have \( b(\xi) \eqm \t_{\cB}(\xi) \).
In words: a sequence is nonrandom with respect to all Bernoulli measures if and
only if it is rejected by a universal combinatorial Bernoulli test.
 \end{theorem}

\begin{formalization}
Planned. The estimate~\eqref{eq:Bernoulli-estim-1} and $(N/(N-n))^{n}\le e^{2}$ are elementary real analysis.
\end{formalization}
 \begin{proof}
  \proves{thm:combinat-Bernoulli}
  \uses{def:g4.2.4, prop:g4.2.5, prop:g4.2.13, propo:Bernoulli-extend}
We first show \( b(\xi)\le \t_{\cB}(\xi) \).
Moreover, we show that \( b(x) \) is an extended
class test for the class of Bernoulli measures.
We only need to check the sum condition, namely that for all
\( 0\le p \le 1 \), and all \( n>0 \) the inequality
\( \sum_{x\in\dB^{n}}B_{p}(x)b(x) \le 1 \) holds.
Indeed, we have
 \begin{align*}
 \sum_{x\in\dB^{n}}B_{p}(x)f(x)
     &= \sum_{k=0}^{n}p^{k}(1-p)^{n-k}\sum_{x\in\dB(n,k)}f(x)
\\  &\le \sum_{k=0}^{n}p^{k}(1-p)^{n-k}\binom{n}{k} = 1.
 \end{align*}
On the other hand, let \( f(x)=\t_{\cB}(x) \), \( x\in\dB^{*} \)
be the extended test for \( \t_{\cB}(\xi) \).
For all integers \( N>0 \) let \( n = \flo{\sqrt{N/2}} \).
Then as \( N \) runs through the integers, \( n \) also runs through all integers.
For \( x\in\dB^{N} \) let \( F(x)=f(x^{\le n}) \).
Since \( f(x) \) is lower semicomputable and monotonic with respect to the prefix
relation, this is also true of \( F(x) \).

We need to estimate \( \sum_{x\in\dB(N,K)} F(x) \).
For this, note that for \( y\in\dB(n,k) \) we have
 \begin{align}\label{eq:hyper}
   \card{\setof{x\in\dB(N,K):y\prefix x}} = \binom{N-n}{K-k}.
 \end{align}
Now for \( 0\le K\le N \) we have
 \begin{equation}\label{eq:Bernoulli-estim}
 \begin{aligned}
      \sum_{x\in\dB(N,K)} F(x) &= 
     \sum_{y\in\dB^{n}}f(y)\card{\setof{x\in\dB(N,K):y\prefix x}}
\\ &= \sum_{k=0}^{n}\binom{N-n}{K-k}\sum_{y\in\dB(n,k)}f(y).
 \end{aligned}\end{equation}
Let us estimate \( \binom{N-n}{K-k}\big/\binom{N}{K} \).
If \( K=k=0 \) then this is 1.
If \( k=n \) then it is \( \frac{K\dotsm(K-n+1)}{N\dotsm(N-n+1)} \).
Otherwise, using \( p=K/N \):
\begin{align}\nonumber
\frac{\binom{N-n}{K-k}}{\binom{N}{K}}
 &= \frac{(N-n)(N-n-1)\dotsm(N-K-(n-k)+1)/(K-k)!}{N(N-1)\dotsm(N-K+1)/K!}
\\\nonumber  &= \frac{K\dotsm(K-k+1)\cdot(N-K)\dotsm(N-K-(n-k)+1)}{N\dotsm(N-n+1)}.
\\\label{eq:Bernoulli-estim-1}
  &\le \frac{K^{k}(N-K)^{n-k}}{(N-n)^{n}} = p^{k}(1-p)^{n-k}\Paren{\frac{N}{N-n}}^{n}.
 \end{align}
Thus in all cases the estimate
 \begin{align*}
\binom{N-n}{K-n}\bigg/\binom{N}{K} \le p^{k}(1-p)^{n-k}\Paren{\frac{N}{N-n}}^{n}
 \end{align*}
holds.
We have
 \begin{align*}
\Paren{\frac{N}{N-n}}^{n} =\Paren{1+\frac{n}{N-n}}^{n}
  \le e^{\frac{N}{2(N-n)}} \le e^{2},
 \end{align*}
since we assumed \( 2n^{2}\le N \).
Substituting into~\eqref{eq:Bernoulli-estim} gives
 \begin{align*}
  \sum_{x\in\dB(N,K)} F(x)
  &\le e^{2}\sum_{k=0}^{n}p^{k}(1-p)^{n-k}\sum_{y\in\dB(n,k)}f(y) \le e^{2},
 \end{align*}
since \( f(x) \) is an extended class test for Bernoulli measures.
It follows that \( e^{-2}F(x)\lem b(x) \), hence also \( F(x) \lem b(x) \).
But we have \( \t_{\cC}(\xi)=\sup_{n}f(\xi^{\le n}) = \sup_{n}F(\xi^{\le n}) \), hence
\( \t_{\cC}(\xi) \lem b(\xi) \).
 \end{proof}


\section{Neutral measure}\label{sec:neutral}

Let \( \t_{\mu}(x) \) be our universal uniform randomness test.
We call a measure \( M \) \df{neutral} if
\( \t_{M}(x) \le 1 \) for all \( x \).
If \( M \) is neutral then no experimental outcome \( x \) could 
refute the theory (hypothesis, model)
that \( M \) is the underlying measure to our experiments.
It can be used as ``apriori probability'', in a Bayesian approach to
statistics.
Levin's theorem says the following:

% Gacs source line 5729
 \begin{theorem}\label{thm:neutral-meas}
  \notready
  \uses{def:g4.1.5}
If the space \( \bX \) is compact then there is a neutral measure over \( \bX \).
 \end{theorem}
 The proof relies on a nontrivial combinatorial fact,
Sperner's Lemma, which also underlies the proof of the Brouwer fixpoint
theorem.
Here is a version of Sperner's Lemma, spelled out in continuous form:


\begin{formalization}
Not planned. The proof needs the compactness of the space of measures over a compact space
(Proposition A.2.50) and Sperner's lemma (Proposition~\ref{propo:Sperner}); both are outside the
agreed scope.
\end{formalization}
% Gacs source line 5737
 \begin{proposition}\label{propo:Sperner}
  \notready
Let \( p_{1},\dots,p_{k} \) be points of some finite-dimensional space
\( \dR^{n} \).
Suppose that there are closed sets \( F_{1},\dots,F_{k} \)
with the property that for every subset \( 1 \le i_{1} < \dots < i_{j} \le k \)
of the indices, the simplex \( S(p_{i_{1}},\dots,p_{i_{j}}) \) spanned by
\( p_{i_{1}},\dots,p_{i_{j}} \) is covered by
the union \( F_{i_{1}} \cup \dots \cup F_{i_{j}} \).
Then the intersection \( \bigcap_{i} F_{i} \) of all these sets is not empty.
 \end{proposition}

The following lemma will also be needed.

% Gacs source line 5750
 \begin{lemma}\label{lem:concentr}
  \uses{def:g4.1.5}
For every closed set \( A \sbs \bX \) and measure \( \mu \), if \( \mu(A)=1 \) then
there is a point \( x\in A \) with \( \t_{\mu}(x) \le 1 \).
 \end{lemma}

\begin{formalization}
Planned over $\dS$ (any set $A$) and over $\dB^{\dN}$ ($A$ closed: the complement of an open set
given by generators, as in Chapter 2's \texttt{openOf}). The last step needs strictness: if
$\t_{\mu}>1$ on $A$ and $\mu(A) = 1$, then $\mu\t_{\mu}>1$.
\end{formalization}
 \begin{proof}
  \proves{lem:concentr}
This follows easily from 
\( \mu\, t_{\mu} = \mu^{x} 1_{A}(x)t_{\mu}(x) \le 1 \).
 \end{proof} 

 \begin{proof}[Proof of Theorem~\protect\ref{thm:neutral-meas}]
  \proves{thm:neutral-meas}
  \notready
  \uses{propo:Sperner, lem:concentr}
For every point \( x \in \bX \), let \( F_{x} \) be the set of measures for which
\( \t_{\mu}(x) \le 1 \).
If we show that for every finite set of points \( x_{1},\dots,x_{k} \), we
have 
 \begin{equation}\label{eq:finite-inters}
  F_{x_{1}}\cap\dots\cap F_{x_{k}} \ne \emptyset,
 \end{equation}
then we will be done.
Indeed, according to Proposition~A.2.50, the compactness
of \( \bX \) implies the compactness of the space \( \bM(\bX) \) of measures.
Therefore if every finite subset of the family \( \setof{F_{x} : x \in \bX} \)
of closed sets has a nonempty intersection, then the whole family has a
nonempty intersection: this intersection consists of the neutral
measures.

To show~\eqref{eq:finite-inters}, let \( S(x_{1},\dots,x_{k}) \) be the set of
probability measures concentrated on \( x_{1},\dots,x_{k} \).
Lemma~\ref{lem:concentr} implies that each such measure belongs to one of
the sets \( F_{x_{i}} \).
Hence \( S(x_{1},\dots,x_{k}) \sbs F_{x_{1}} \cup \dots \cup F_{x_{k}} \),
and the same holds for every subset of the indices \( \{1,\dots,k\} \).
Sperner's Lemma~\ref{propo:Sperner} implies 
\( F_{x_{1}} \cap \dots \cap F_{x_{k}} \ne\emptyset \).
 \end{proof}

When the space is not compact, there are generally no neutral probability
measures, as shown by the following example.

% Gacs source line 5788
 \begin{proposition}\label{thm:no-neutral}
  \uses{def:g4.1.5}
Over the discrete space \( \bX = \dN \) of natural numbers,
there is no neutral measure.
 \end{proposition}

\begin{formalization}
Planned, with $\dN$ identified with $\dS$ by the standard enumeration. The test~\eqref{eq:no-neutral} is lower
semicomputable in $(x,\mu)$ in the sense of Definition~\ref{our:lsc}, because $\sum_{y<x}\mu(y)$
is a finite sum and the condition on it is strict.
\end{formalization}
 \begin{proof}
  \proves{thm:no-neutral}
  \uses{our:lsc}
It is sufficient to 
construct a randomness test \( t_{\mu}(x) \) with the property that for
every measure \( \mu \), we have \( \sup_{x} t_{\mu}(x) = \infty \).
Let 
 \begin{equation}\label{eq:no-neutral}
   t_{\mu}(x) = \sup\setof{ k \in \dN: \sum_{y<x}\mu(y) > 1-2^{-k}}.
 \end{equation}
By its construction, this is a lower semicomputable function with
\( \sup_{x} t_{\mu}(x) = \infty \).
It is a test if \( \sum_{x}\mu(x)t_{\mu}(x) \le 1 \).
We have
 \[
  \sum_{x} \mu(x) t_{\mu}(x) 
= \sum_{k>0} \sum_{t_{\mu}(x) \ge k} \mu(x)
< \sum_{k>0} 2^{-k} \le 1.
 \]
 \end{proof}

Using a similar construction over the space \( \dN^{\dN} \) of infinite
sequences of natural numbers, we could show that 
for every measure \( \mu \) there is a sequence \( x \) with \( \t_{\mu}(x)=\infty \).

Proposition~\ref{thm:no-neutral} is a little misleading, since \( \dN \) can be
compactified into \( \ol\dN = \dN \cup \{\infty\} \)
(as in Part~1 of Example~A.1.22).
Theorem~\ref{thm:neutral-meas} implies that there is a neutral probability
measure \( M \) over the compactified space \( \ol\dN \).
Its restriction to \( \dN \) is, of course, not a probability measure, since it
satisfies only \( \sum_{x < \infty} M(x) \le 1 \).
We called these functions \df{semimeasures}.

 \begin{remark}\label{rem:compactify}\
 \begin{enumerate}

   \item
 It is easy to see that 
Theorem~\ref{thm:test-charac-discr} characterizing randomness in terms of
complexity holds also for the space \( \ol\dN \).

   \item
The topological space of semimeasures
over \( \dN \) is not compact, and there is no neutral one among them.
Its topology is not the same as what we get when we restrict the topology
of probability measures over \( \ol\dN \) to \( \dN \).
The difference is that over \( \dN \), for example the set of measures
\( \setof{\mu : \mu(\dN) > 1/2} \) is closed, since \( \dN \) (as the whole space)
is a closed set.
But over \( \ol\dN \), this set is not necessarily closed, since
\( \dN \) is not a closed subset of \( \ol\dN \).
 \end{enumerate}
 \end{remark}

Neutral measures are not too simple, even over \( \ol\dN \), as the following
theorem shows.

% Gacs source line 5848
 \begin{theorem}\label{thm:no-upper-semi-neutral}
  \notready
  \uses{thm:test-charac-discr}
There is no neutral measure over \( \ol\dN \) that is upper
semicomputable over \( \dN \) or lower semicomputable over \( \dN \).
 \end{theorem}

\begin{formalization}
Not planned: it is about neutral measures over the compactified space $\ol\dN$, outside the
specialization.
\end{formalization}
 \begin{proof}
  \proves{thm:no-upper-semi-neutral}
  \notready
Let us assume that \( \nu \) is a measure that is upper semicomputable over
\( \dN \).
Then the set
 \[
   \setof{(x,r) : x \in\dN,\; r\in\dQ,\; \nu(x) < r}
 \]
is recursively enumerable: let \( (x_{i},r_{i}) \) be a particular
enumeration. 
For each \( n \), let \( i(n) \) be the first \( i \) with \( r_{i} < 2^{-n} \), and let
\( y_{n} = x_{i(n)} \).
Then \( \nu(y_{n}) < 2^{-n} \), and at the same time \( K(y_{n}) \lea K(n) \).
As mentioned, in Remark~\ref{rem:compactify}, 
Theorem~\ref{thm:test-charac-discr} characterizing randomness in terms of
complexity holds also for the space \( \ol\dN \).
Thus, 
 \[
  \d_{\nu}(y_{n}) \eqa -\log\nu(y_{n}) - K(y_{n} \mvert \nu) \gea n - K(n).
 \]
Suppose now that \( \nu \) is lower semicomputable over \( \dN \).
The proof for this case is longer.
We know that \( \nu \) is the monotonic limit of a recursive sequence
\( i\mapsto \nu_{i}(x) \) of recursive semimeasures with rational values
\( \nu_{i}(x) \).
For every \( k=0,\dots,2^{n}-2 \), let 
 \begin{align*}
  V_{n,k} &= \setof{\mu \in \cM(\ol\dN) : 
k\cdot 2^{-n} < \mu(\{0,\dots,2^{n}-1\}) < (k+2)\cdot 2^{-n}},
\\       J &= \setof{(n,k): k\cdot 2^{-n} < \nu(\{0,\dots,2^{n}-1\})}.
 \end{align*}
The set \( J \) is recursively enumerable.
Let us define the functions \( j:J\to\dN \) and \( x:J\to\{0,\dots,2^{n}-1\} \)
as follows: \( j(n,k) \) is the smallest \( i \) with 
\( \nu_{i}(\{0,\dots,2^{n}-1\}) > k\cdot 2^{-n} \), and
 \[
   x_{n,k} = \min\setof{y < 2^{n}: \nu_{j(n,k)}(y) < 2^{-n+1}}.
 \]
Let us define the function \( f_{\mu}(x,n,k) \) as follows.
We set \( f_{\mu}(x,n,k)=2^{n-2} \) if the following conditions hold:
 \begin{alphenum}
  \item\label{i:no-upper-semi-neutral.mu-global} \( \mu \in V_{n,k} \);
  \item\label{i:no-upper-semi-neutral.mu-upper} \( \mu(x) < 2^{-n+2} \);
  \item\label{i:no-upper-semi-neutral.unique}
 \( (n,k) \in J \) and \( x=x_{n,k} \).
 \end{alphenum}
Otherwise, \( f_{\mu}(x,n,k)=0 \).
Clearly, the function \( (\mu,x,n,k) \mapsto f_{\mu}(x,n,k) \) is lower
semicomputable.
Condition~\eqref{i:no-upper-semi-neutral.mu-upper} implies
 \begin{equation}\label{eq:no-upper-semi-neutral.n-k-test}
 \sum_{y} \mu(y) f_{\mu}(y,n,k) \le
    \mu(x_{n,k})f_{\mu}(x_{n,k},n,k) < 2^{-n+2}\cdot 2^{n-2} = 1.
 \end{equation}
Let us show that \( \nu \in V_{n,k} \) implies
 \begin{equation}\label{eq:found-bad}
 f_{\nu}(x_{n,k},n,k) = 2^{n-2}.
 \end{equation}
Consider \( x=x_{n,k} \).
Conditions~\eqref{i:no-upper-semi-neutral.mu-global} 
and~\eqref{i:no-upper-semi-neutral.unique} are satisfied by definition.
Let us show that condition~\eqref{i:no-upper-semi-neutral.mu-upper} is also 
satisfied.
Let \( j=j(n,k) \).
By definition, we have \( \nu_{j}(x) < 2^{-n+1} \).
Since by definition \( \nu_{j}\in V_{n,k} \) and \( \nu_{j} \le \nu \in V_{n,k} \),
we have 
 \[
  \nu(x) \le \nu_{j}(x) + 2^{-n+1} < 2^{-n+1} + 2^{-n+1} = 2^{-n+2}.
 \]
Since all three conditions~\eqref{i:no-upper-semi-neutral.mu-global},
\eqref{i:no-upper-semi-neutral.mu-upper}
and~\eqref{i:no-upper-semi-neutral.unique} are satisfied, we have 
shown~\eqref{eq:found-bad}.
Now we define
 \[
   g_{\mu}(x) = \sum_{n\ge 2}\frac{1}{n(n+1)}\sum_{k}f_{\mu}(x,n,k).
 \]
Let us prove that \( g_{\mu}(x) \) is a uniform test.
It is lower semicomputable by definition, so we only need to prove
\( \sum_{x} \mu(x) f_{\mu}(x) \le 1 \).
For this, let \( I_{n,\mu} = \setof{k: \mu\in V_{n,k}} \).
Clearly by definition, \( \card{I_{n,\mu}}\le 2 \).
We have, using this last fact and the test 
property~\eqref{eq:no-upper-semi-neutral.n-k-test}:
 \[
  \sum_{x} \mu(x) g_{\mu}(x) = 
  \sum_{n\ge 2}\frac{1}{n(n+1)}
  \sum_{k\in I_{n,\mu}} \sum_{x}\mu(x) f_{\mu}(x,n,k)
       \le  \sum_{n\ge 2}\frac{1}{n(n+1)}\cdot 2 \le 1.
 \]
Thus, \( g_{\mu}(x) \) is a uniform test.
If \( \nu\in V_{n,k} \) then we have
 \[
 \t_{\nu}(x_{n,k}) 
\gem g_{\nu}(x_{n,k}) \ge \frac{1}{n(n+1)}f_{\mu}(x_{n,k},n,k) \ge 
  \frac{2^{n-2}}{n(n+1)}.
 \]
Hence \( \nu \) is not neutral.
 \end{proof}

% Gacs source line 5952
 \begin{remark}
  \label{rem:g4.3.5}
In~\cite{LevinUnif76} and~\cite{LevinRandCons84},
Levin imposed extra conditions on tests which allow to find a lower
semicomputable neutral semimeasure.
 \end{remark}

The universal lower semicomputable
semimeasure \( \m(x) \) has a certain property similar to neutrality.
According to Theorem~\ref{thm:defic-charac-cpt-seqs} specialized to
one-element sequences, for every computable measure \( \mu \) we have
\( \d_{\mu}(x) \eqa -\log\mu(x) - K(x) \)
(where the constant in \( \eqa \) depends on \( \mu \)).
So, for computable measures, the expression
 \begin{equation}\label{eq:ol-d}
  \ol\d_{\mu}(x) = -\log\mu(x) - K(x)
 \end{equation}
can serve as a reasonable deficiency of randomness.
(We will also use the test \( \ol\t = 2^{\ol\d} \).)
If we substitute \( \m \) for \( \mu \) in \( \ol\d_{\mu}(x) \), we get 0.
This substitution is not justified, of course.
The fact that \( \m \) is not a probability 
measure can be helped, using compactification as above, 
and extending the notion of randomness tests. 
But the test \( \ol\d_{\mu} \) can replace \( \d_{\mu} \)
only for computable \( \mu \), while \( \m \) is not computable.
Anyway, this is the sense in which all outcomes might
be considered random with respect to \( \m \), and the heuristic
sense in which \( \m \) may be considered ``neutral''.

\section{Monotonicity, quasi-convexity/concavity}

 Some people find that \( \mu \)-tests as defined in Definition~\ref{def:mu-test} are
 too general, in case \( \mu \) is a non-computable measure.
 In particular, randomness with respect to computable measures has 
 a certain---intuitively meaningful---monotonicity property:
 roughly, if \( \nu \) is greater than \( \mu \) then if
 \( x \) is random with respect to \( \mu \), it should also be random with respect to
 \( \nu \).

% Gacs source line 5991
  \begin{proposition}\label{propo:test-computable-mon}
  \uses{def:g4.1.5}
 For computable measures \( \mu,\nu \) we have for all rational \( c>0 \):
  \begin{align}\label{eq:test-computable-mon}
   2^{-k}\mu\le \nu \imp \d_{\nu}(x)\lea \d_{\mu}(x)+ k + K(k).
  \end{align}
 Here the constant in \( \lea \) depends on \( \mu,\nu \), but not on \( k \).
  \end{proposition}

\begin{formalization}
Planned for computable $\mu,\nu$ over either space, with the constant independent of $k$. The sum
is built with Theorem~\ref{thm:trim}, trimming $2^{-k-1}\t_{\nu}$ to $\mu$-integral at most 1,
which leaves it unchanged whenever $2^{-k}\mu\le\nu$; the weights $2^{-K(k)}$ sum to at most 1 by
Kraft's inequality.
\end{formalization}
 \begin{proof}
  \proves{propo:test-computable-mon}
  \uses{thm:trim, l:kraft}
 We have \( 1\ge\nu\t_{\nu}\ge 2^{-k}\mu\t_{\nu} \),
 hence \( 2^{-k}\t_{\nu} \) is a \( \mu \)-test.
 Using the method of Theorem~\ref{thm:trim} in finding universal tests, 
 one can show that the sum
  \begin{align*}
    \sum_{k: 2^{-k}\mu\t_{\nu}<1} 2^{-k-K(k)}\t_{\nu}
  \end{align*}
 is a \( \mu \)-test, and hence \( \lem \t_{\mu} \).
 Therefore this is true of each member of the sum, which is just what the theorem
 claims.
 \end{proof}

 There are other properties true for tests on computable measures that we may
 want to require for all measures.
 For the following properties, let us define quasi-convexity, which is a
 weakening of the notion of convexity.

% Gacs source line 6016
 \begin{definition}\label{def:quasi-convex}
 A function \( f:V\to\dR \) defined on a vector space \( V \) is called
 \df{quasi-convex} if for every real number \( x \) the set \( \setof{v: f(v)\le x} \)
 is convex.
 It is  \df{quasi-concave} if \( -f \) is quasi-convex.
 \end{definition}

\begin{formalization}
Planned as Mathlib's \texttt{QuasiconvexOn} and \texttt{QuasiconcaveOn}; nothing new to define.
\end{formalization}

 It is easy to see that quasi-convexity is equivalent to the inequality
  \begin{align*}
  f(\lg u + (1-\lg)v)\le f(u)\lor f(v)
  \end{align*}
 for all \( u,v \) and \( 0<\lg<1 \), while quasi-concavity is equivalent to
  \begin{align*}
  f(\lg u + (1-\lg)v)\ge f(u)\land f(v)
  \end{align*}

The uniform test with respect to computable measures is approximately 
both quasi-convex and quasi-concave.
Let \( \nu=\lg\mu_{1}+(1-\lg)\mu_{2} \).

Quasi-convexity means, roughly, that if \( x \) is random with respect to both
\( \mu_{1} \) and \( \mu_{2} \) then it is also random with respect to \( \nu \).
This property strengthens monotonicity in the cases where it applies.

% Gacs source line 6040
  \begin{proposition}\label{propo:test-computable-conv}
  \uses{def:g4.1.5, def:quasi-convex}
 Let \( \mu_{1},\mu_{2} \) be computable measures and \( 0<\lg<1 \) computable,
with \( \nu=\lg\mu_{1}+(1-\lg)\mu_{2} \).
 Then we have
  \begin{align*}
   \d_{\nu}(x) \lea\d_{\mu_{1}}(x)\lor\d_{\mu_{2}}(x) + K(\lg).
  \end{align*}
  \end{proposition}

\begin{formalization}
Planned with a constant depending on $\mu_{1},\mu_{2},\lg$. The explicit term $K(\lg)$ is a
refinement through descriptions, as for the constants $K(d)+K(P)$ in Chapter 2
(\texttt{IsLSCDescription}).
\end{formalization}
 \begin{proof}
  \proves{propo:test-computable-conv}
  \uses{coroll:trim}
The relation \( 1\ge\nu\t_{\nu}=\lg\mu_{1}\t_{\nu}+(1-\lg)\mu_{2}\t_{\nu} \)
implies \( 1\ge \mu_{i}\t_{\nu} \) for some \( i\in\{1,2\} \).
Then \( \d_{\nu}\lea\d_{\mu_{i}}+K(\lg) \) (since \( \lg \) was used to define \( \nu \) and
thus \( \t_{\nu} \)).
 \end{proof}

Quasi-concavity means, roughly, that if \( x \) is non-random with respect to both
\( \mu_{1} \) and \( \mu_{2} \) then it is also nonrandom with respect to \( \nu \):

% Gacs source line 6058
  \begin{proposition}\label{propo:test-computable-conc}
  \uses{def:g4.1.5, def:quasi-convex}
 Let \( \mu_{1},\mu_{2} \) be computable measures and \( 0<\lg<1 \) arbitrary
(not even necessarily computable), with \( \nu=\lg\mu_{1}+(1-\lg)\mu_{2} \).
 Then we have
  \begin{align*}
   \d_{\nu} \gea\d_{\mu_{1}}\land\d_{\mu_{2}}.
  \end{align*}
  \end{proposition}

\begin{formalization}
Planned. Since $\lg$ is arbitrary, $\nu$ is in general not computable and the statement needs the
uniform test. The constant depends only on $\mu_{1},\mu_{2}$: it comes from
Corollary~\ref{coroll:trim} applied to the fixed function $\t_{\mu_{1}}\land\t_{\mu_{2}}$.
\end{formalization}
  \begin{proof}
  \proves{propo:test-computable-conc}
  \uses{coroll:trim}
 The function \( \t_{\mu_{1}}\land\t_{\mu_{2}} \) 
is lower semicomputable, and is a \( \mu_{i} \)-test for each \( i \).
Therefore it is also a \( \nu \)-test, and as such is \( \lem\t_{\nu} \).
Here, the constant in the \( \lem \) depends only on (the programs for)
\( \mu_{1},\mu_{2} \), and not on \( \lg \).
  \end{proof}

 These properties do not survive for arbitrary measures and arbitrary constants.

% Gacs source line 6076
  \begin{example}\label{example:monotonicity-counter}\
 Let measure \( \mu_{1} \) be uniform over the interval \( \clint{0}{1/2} \), 
let \( \mu_{2} \) be uniform over \( \clint{1/2}{1} \).
For \( 0\le x \le 1 \) let
 \begin{align*}
      \nu_{x} &= (1-x)\mu_{1} + x\mu_{2}.
  \end{align*}
Then \( \nu_{1/2} \) is uniform over \( \clint{0}{1} \).
Let \( \phi_{x} \) be the uniform distribution over \( \clint{x}{x+1/2} \),
and \( \psi_{x} \) the uniform distribution over
\( \clint{0}{x}\cup\clint{x+1/2}{1} \).

Let \( p<1/2 \) be random with respect to the uniform
distribution \( \nu_{1/2} \).

 \begin{enumerate}
  \item The relations
 \begin{align*}
  \nu_{1/2}\le p^{-1}\nu_{p},
  \quad \d_{\nu_{1/2}}(p)<\infty,\quad\d_{\nu_{p}}(p)=\infty
 \end{align*} 
show that any statement analogous to the monotonicity property of
Proposition~\ref{propo:test-computable-mon}
fails when the measures involved are not required to be computable.

  \item
For rational \( r \) with \( 0<r<p \) the relations
\begin{align*}
   \nu:=(1-p)\nu_{r} + p\nu_{1-r},
\quad \d_{\nu}(p)=\infty,
\quad\d_{\nu_{r}}(p)<\infty,
\quad \d_{\nu_{1-r}}(p)<\infty 
 \end{align*}
provide a similar counterexample to Proposition~\ref{propo:test-computable-conv}.

  \item
The relations
 \begin{align*}
 \nu_{1/2} = (\nu_{p}+\nu_{1-p})/2, 
 \quad \d_{\nu_{1/2}}(p)<\infty,\quad \d_{\nu_{p}}(p)=\d_{\nu_{1-p}}(p)=\infty
 \end{align*}
provide a similar counterexample to Proposition~\ref{propo:test-computable-conc}.

The following counterexample relies on a less subtle effect:
 \begin{align*}
\nu_{1/2}=(\phi_{p}+\psi_{p})/2,
\quad\d_{\nu_{1/2}}(p)<\infty,
\quad\d_{\phi_{p}}(p)=\d_{\psi_{p}}(p)=\infty,
 \end{align*}
since
as a boundary point of the support, \( p \) is computable from both \( \phi_{p} \) and
\( \psi_{p} \) in a uniform way.
 \end{enumerate}
For a complete proof, uniform tests must be provided for each of the cases: this
is left as exercise for the reader.
  \end{example}

\begin{formalization}
Not planned: an example over $\clint{0}{1}$, whose proofs are left as an exercise.
\end{formalization}

The non-monotonicity example could be used to argue that the 
we allowed too many \( \mu \)-tests, that the test \( \t_{\mu}(x) \)
should not be allowed to depend on properties of \( \mu \) that exploit the
computational properties of \( \mu \) so much stronger than its quantitative
properties.

\section{Algorithmic entropy}

Some properties of description complexity make it a good
expression of the idea of individual information content.

\subsection{Entropy}
The entropy of a discrete probability distribution $\mu$ is defined as
 \[
   \cH(\mu) = - \sum_{x} \mu(x) \log \mu(x).
 \]
To generalize entropy to continuous distributions the
\df{relative entropy} is defined as follows.
Let $\mu,\nu$ be two measures, where $\mu$ is taken (typically, but not
always), to be a probability measure, and $\nu$ another measure, that can
also be a probability measure but is most frequently not.
We define the \df{relative entropy} $\cH_{\nu}(\mu)$ as follows.
If $\mu$ is not absolutely continuous with respect to $\nu$ then 
$\cH_{\nu}(\mu) = -\infty$.
Otherwise, writing 
 \[
  \frac{d\mu}{d\nu} = \frac{\mu(dx)}{\nu(dx)} =: f(x)
 \] 
for the (Radon-Nikodym) derivative 
(density) of $\mu$ with respect to $\nu$, we define
 \[
   \cH_{\nu}(\mu) = - \int \log\frac{d\mu}{d\nu} d\mu 
     = - \mu^{x} \log\frac{\mu(dx)}{\nu(dx)} = -\nu^{x} f(x) \log f(x).
 \]
Thus, $\cH(\mu) = \cH_{\#}(\mu)$ is a special case.

% Gacs source line 6169
 \begin{example}
  \label{ex:g4.5.1}
 Let $f(x)$ be a probability density function for the distribution
$\mu$ over the real line, and let $\lg$ be the Lebesgue measure there.
Then
 \[
   \cH_{\lg}(\mu) = -\int f(x) \log f(x) d x.
 \]
 \end{example}

In information theory and statistics, when both $\mu$ and $\nu$ are
probability measures, then $-\cH_{\nu}(\mu)$ is also denoted
$D(\mu \parallel \nu)$, and called (after Kullback) the
information divergence of the two measures.
It is frequently used in the role of a distance between $\mu$ and $\nu$.
It is not symmetric, but can be shown to obey the triangle inequality, and
to be nonnegative.
Let us prove the latter property: in our terms, it says that relative
entropy is nonpositive when both $\mu$ and $\nu$ are probability measures.

% Gacs source line 6188
 \begin{proposition}\label{propo:Kullback-pos}
  \uses{def:g3.1.1}
Over a space $X$, we have
 \begin{equation}\label{eq:Kullback-pos}
   \cH_{\nu}(\mu) \le -\mu(X) \log\frac{\mu(X)}{\nu(X)}.
 \end{equation}
 In particular, if $\mu(X) \ge \nu(X)$ then $\cH_{\nu}(\mu) \le 0$.
 \end{proposition}

\begin{formalization}
Planned over $\dS$, where the density is $\mu(x)/\nu(x)$, in the conventions of Chapter 3's
classical entropy (\texttt{Classical/Entropy.lean}): it is Gibbs' inequality
(\texttt{entropy\_le\_crossEntropy}, \texttt{log\_sum\_inequality}). Densities over
$\dB^{\dN}$ are not planned.
\end{formalization}
 \begin{proof}
  \proves{propo:Kullback-pos}
  \uses{thm:CsiszarConc}
The inequality $- a \ln a \le -a\ln b + b-a$
expresses the concavity of the logarithm function.
Substituting $a = f(x)$ and $b = \mu(X)/\nu(X)$ 
and integrating by $\nu$:
 \begin{align*}
 (\ln 2) \cH_{\nu}(\mu) &= 
 -\nu^{x} f(x) \ln f(x) \le -\mu(X) \ln\frac{\mu(X)}{\nu(X)}
  + \frac{\mu(X)}{\nu(X)} \nu(X) - \mu(X)
\\ &= -\mu(X) \ln\frac{\mu(X)}{\nu(X)},
 \end{align*}
giving~\eqref{eq:Kullback-pos}.
 \end{proof}

\subsection{Algorithmic entropy}

Le us recall some facts on description complexity.
Let us fix some (finite or infinite) 
alphabet $\Sg$ and consider the discrete space $\Sg^{*}$.

The universal lower semicomputable
semimeasure $\m(x)$ over $\Sg^{*}$ was defined 
in Definition~\ref{def:m}.
It is possible to turn $\m(x)$ into a measure, by compactifying the
discrete space $\Sg^{*}$ into 
 \[
  \ol{\Sg^{*}}=\Sg^{*}\cup\{\infty\}
 \]
(as in part~1 of Example~A.1.22;
this process makes sense also for a constructive discrete space),
and setting $\m(\infty) = 1-\sum_{x\in\Sg^{*}} \m(x)$.
The extended measure $\m$ is not quite lower semicomputable since
the number $\mu(\ol{\Sg^{*}} \xcpt \{0\})$ is not necessarily
lower semicomputable.

% Gacs source line 6230
 \begin{remark}
  \label{rem:g4.5.3}
A measure $\mu$ is computable over $\ol{\Sg^{*}}$ if and only if the
function $x \mapsto \mu(x)$ is computable for $x \in \Sg^{*}$.
This property does not imply that the number
 \[
 1 - \mu(\infty) = \mu(\Sg^{*}) = \sum_{x\in\Sg^{*}} \mu(x)
 \]
 is computable.
 \end{remark}

Let us allow, for a moment, measures $\mu$ that are not probability
measures: they may not even be finite.
Metric and computability can be extended to this case, the universal test
$\t_{\mu}(x)$ can also be generalized.
The Coding Theorem~\ref{thm:coding}
and other considerations suggest the introduction of
the following notation, for an arbitrary measure $\mu$:

% Gacs source line 6248
 \begin{definition}
  \label{def:g4.5.4}
  \uses{def:g4.1.5, our:m-cond, thm:coding}
We define the \df{algorithmic entropy} of a point $x$ with respect to measure
$\mu$ as
 \begin{equation}\label{eq:alg-ent}
   H_{\mu}(x) = -\d_{\mu}(x) = -\log\t_{\mu}(x).
 \end{equation}
 \end{definition}

\begin{formalization}
Planned, multiplicatively: $2^{-H_{\mu}(x)} = \t_{\mu}(x)$. Here $\mu$ ranges over
Definition~\ref{our:measure} and the counting measure $\#$ over $\dS$. Then $\t_{\#} = \m$ is
Chapter 1's universal semimeasure, so $K(x)\eqa H_{\#}(x)$ is the coding theorem.
\end{formalization}
Then, with $\#$ defined as the counting measure over the discrete set
$\Sg^{*}$ (that is, $\#(S) = \card{S}$), we have
 \[
   K(x) \eqa H_{\#}(x).
 \]
This allows viewing $H_{\mu}(x)$ as a generalization of description
complexity.

The following theorem generalizes an earlier known theorem stating that
over a discrete space, for a computable measure,
entropy is within an additive constant the same as ``average complexity'':
$\cH(\mu) \eqa \mu^{x} K(x)$.

% Gacs source line 6268
 \begin{theorem}
  \label{thm:g4.5.1}
  \uses{def:g4.5.4, thm:g3.1.3}
Let $\mu$ be a probability measure.
Then we have
 \begin{equation}\label{eq:entropy-less-avg-algentr}
  \cH_{\nu}(\mu) \le \mu^{x} H_{\nu}(x \mvert \mu).
 \end{equation}
If $X$ is a discrete space then the following estimate also holds:
  \begin{equation}\label{eq:entropy-gea-avg-algentr}
  \cH_{\nu}(\mu) \gea \mu^{x} H_{\nu}(x \mvert \mu).
  \end{equation}
 \end{theorem}

\begin{formalization}
Planned for $X = \dS$ only, with $\nu$ a measure of Definition~\ref{our:measure} or $\#$; over
other spaces the first inequality needs densities. For $\nu = \#$ it reads
$\cH(\mu)\le\sum_{x}\mu(x)K(x \mvert \mu)\lea\cH(\mu)$, which generalizes Theorem~3.1.3
(\texttt{Classical.entropy\_expectedComplexity}) to noncomputable $\mu$, with
$K(x \mvert \mu)$ in place of $K(x)$ and no $K(\mu)$ term.
\end{formalization}
 \begin{proof}
  \proves{thm:g4.5.1}
  \uses{propo:Kullback-pos, our:lsc-props}
Let
$\dg$ be the measure with density $\t_{\nu}(x \mvert \mu)$ with respect to
$\nu$: $\t_{\nu}(x \mvert \mu) = \frac{\dg(dx)}{\nu(dx)}$.
Then $\dg(X) \le 1$.
It is easy to see from the maximality property of $\t_{\nu}(x \mvert \mu)$
that $\t_{\nu}(x \mvert \mu) > 0$, therefore according to
Proposition~A.2.23, we have
$\frac{\nu(dx)}{\dg(dx)} = \Paren{\frac{\dg(dx)}{\nu(dx)}}^{-1}$.
Using Proposition~A.2.23 and~\ref{propo:Kullback-pos}: 
 \begin{align*}
     \cH_{\nu}(\mu)     &= - \mu^{x} \log\frac{\mu(dx)}{\nu(dx)},
\\   - \mu^{x} H_{\nu}(x \mvert \mu) &=  \mu^{x} \log \frac{\dg(dx)}{\nu(dx)}
              = - \mu^{x} \log \frac{\nu(dx)}{\dg(dx)},
\\    \cH_{\nu}(\mu) - \mu^{x} H_{\nu}(x \mvert \mu)
     &= - \mu^{x} \log \frac{\mu(dx)}{\dg(dx)} 
  \le -\mu(X) \log\frac{\mu(X)}{\dg(X)} \le 0.
 \end{align*}
This proves~\eqref{eq:entropy-less-avg-algentr}.

Over a discrete space $X$, the function
$(x,\mu,\nu) \mapsto \frac{\mu(dx)}{\nu(dx)} = \frac{\mu(x)}{\nu(x)}$
is computable, therefore by the maximality property of
$H_{\nu}(x \mvert \mu)$ we have 
$\frac{\mu(dx)}{\nu(dx)} \lem \t_{\nu}(x \mvert \mu)$,
hence $\cH_{\nu}(\mu) = -\mu^{x} \log \frac{\mu(dx)}{\nu(dx)} 
 \gea \mu^{x} H_{\nu}(x \mvert \mu)$.
 \end{proof} 


\subsection{Addition theorem}
The Addition Theorem~\eqref{eq:Haddit}
can be generalized to the algorithmic entropy $H_{\mu}(x)$
introduced in~\eqref{eq:alg-ent} (a somewhat similar generalization appeared
in~\cite{VovkVyugin93}).
The generalization, defining $H_{\mu,\nu} = H_{\mu\times\nu}$, is
 \begin{equation}\label{eq:addition-general}
   H_{\mu,\nu}(x,y)\eqa
  H_\mu(x \mvert \nu)+ H_\nu(y \mvert x,\; H_\mu(x \mvert \nu),\; \mu).
 \end{equation}
Before proving the general addition theorem, we establish a few useful
facts.

% Gacs source line 6321
 \begin{proposition}\label{propo:int.H.of.xy}
  \uses{def:g4.5.4}
 We have
 \[
  H_{\mu}(x \mvert \nu) \lea -\log \nu^{y} 2^{-H_{\mu,\nu}(x, y)}.
 \]
 \end{proposition}

\begin{formalization}
Planned over $\dS$. The product $\mu\times\nu$ is $\ang{x,y}\mapsto\mu(x)\nu(y)$ on pairs coded by
\texttt{pairCode}, a measure of Definition~\ref{our:measure}, computable from names of $\mu$ and
$\nu$ (Theorem~\ref{our:lsc-props}(d)).
\end{formalization}
 \begin{proof}
  \proves{propo:int.H.of.xy}
  \uses{our:lsc-props}
The function $f(x,\mu,\nu)$ that is the right-hand side, is upper
semicomputable by definition, and obeys $\mu^{x}2^{-f(x,\mu,\nu)} \le 1$.
Therefore the inequality follows from the minimum property of
$H_{\mu}(x)$.
 \end{proof}

Let $z \in \dN$, then the inequality
 \begin{equation}\label{eq:H.x.cond.z}
  H_{\mu}(x) \lea K(z) + H_{\mu}(x \mvert z)
 \end{equation}
 will be a simple consequence of the general addition theorem.
The following lemma, needed in the proof of the theorem,
generalizes this inequality somewhat:

% Gacs source line 6342
  \begin{lemma}\label{lem:H.x.cond.z}
  \uses{def:g4.5.4}
 For a com\-put\-able func\-tion $(y,z) \mapsto f(y,z)$ over $\dN$,
we have
 \[
  H_{\mu}(x \mvert y) \lea K(z) + H_{\mu}(x \mvert f(y,z)).
 \]
 \end{lemma}

\begin{formalization}
Planned. $K(z)$ is Chapter 1's prefix complexity, and $\sum_{z}2^{-K(z)}\le 1$ is Kraft's
inequality.
\end{formalization}
 \begin{proof}
  \proves{lem:H.x.cond.z}
  \uses{our:lsc-props, l:kraft}
  The function
 \[
  (x,y,\mu) \mapsto g_{\mu}(x, y)=\sum_{z} 2^{-H_{\mu}(x \mvert f(y,z))-K(z)}
 \]
 is lower semicomputable,
and $\mu^{x} g_{\mu}(x, y) \le \sum_{z} 2^{-K(z)} \le 1$.
Hence $g_{\mu}(x, y) \lem 2^{-H_{\mu}(x \mvert y)}$.
The left-hand side is a sum, hence the inequality holds for each
element of the sum: just what we had to prove.
 \end{proof}

The following monotonicity property will be needed:

% Gacs source line 6363
\begin{lemma}\label{lem:mon}
  \uses{def:g4.5.4}
 For $i < j$ we have
 \[
   i + H_\mu(x\mvert i) \lea j + H_\mu(x\mvert j) .
 \]
\end{lemma}
\begin{proof}
  \proves{lem:mon}
  \uses{lem:H.x.cond.z}
  From Lemma~\ref{lem:H.x.cond.z}, with $f(i, n)=i + n$ we have
 \[
   H_{\mu}(x \mvert i) - H_{\mu}(x \mvert j) \lea K(j-i) \lea j-i.
 \]
\end{proof}

Let us generalize the minimum property of $H_{\mu}(x)$.

% Gacs source line 6378
 \begin{proposition}\label{propo:univ-test-gener}
  \uses{def:g4.5.4, coroll:trim-univ}
Let $(y,\nu)\mapsto F_{\nu}(y)$ be an upper semicomputable
function with values in $\ol\dZ=\dZ\cup\{-\infty,\infty\}$.
Then by Corollary~\ref{coroll:trim-univ} among the lower semicomputable functions
$(x,y,\nu) \mapsto g_{\nu}(x,y)$ with $\nu^{x}g_{\nu}(x,y)\le 2^{-F_{\nu}(y)}$
there is one that is maximal to within a multiplicative constant.
Choosing $f_{\nu}(x,y)$ as such a function we have
for all $x$ with $F_{\nu}(y) > -\infty$:
 \begin{align*}
  f_{\nu}(x,y) &\eqm 2^{-F_{\nu}(y)}\t_{\nu}(x \mvert y, F_{\nu}(y)),
 \end{align*}
or in logarithmic notation 
$-\log f_{\nu}(x,y) \eqa F_{\nu}(y)+H_{\nu}(x \mvert y, F_{\nu}(y))$.
 \end{proposition}

\begin{formalization}
Planned for $x$ over $\dS$, conditions $y$ as in Definition~\ref{our:m-cond}, and integer-valued
$F$. Substituting the upper semicomputable $F_{\nu}(y)$ for $m$ in $g_{\nu}(x,y,m)$, which is
decreasing in $m$, is Theorem~\ref{our:lsc-props}(d).
\end{formalization}
 \begin{proof}
  \proves{propo:univ-test-gener}
  \uses{thm:trim, lem:mon, our:lsc-props}
  To prove the inequality $\gem$, define
 \[
   g_{\nu}(x,y,m) =\max_{i\ge m}2^{-i}\t_{\nu}(x \mvert y, i).
 \]
  Function $g_{\nu}(x,y,m)$ is lower semicomputable and decreasing
in $m$.
Therefore 
 \[
  g_{\nu}(x,y)  = g_{\nu}(x,y,F_{\nu}(y))
 \]
is also lower semicomputable since it is obtained by substituting an
upper semicomputable function for $m$ in $g_{\nu}(x,y,m)$.
The multiplicative form of Lemma~\ref{lem:mon} implies
 \begin{align*}
       g_{\nu}(x,y,m) &\eqm 2^{-m}\t_{\nu}(x \mvert y, m),
\\   g_{\nu}(x,y) &\eqm 2^{-F_{\nu}(y)}\t_{\nu}(x \mvert y, F_{\nu}(y)).
 \end{align*}
We have, since $\t_{\nu}$ is a test:
 \begin{align*}
   \nu^{x}2^{-m}\t_{\nu}(x \mvert y, m) &\le 2^{-m},
\\   \nu^{x}g_{\nu}(x,y) &\lem 2^{-F_{\nu}(y)},
 \end{align*}
implying $g_{\nu}(x,y)\lem f_{\nu}(x,y)$ by the optimality of $f_{\nu}(x,y)$.

To prove the upper bound, consider all lower semicomputable functions
$\phi_{e}(x,y,m,\nu)$ ($e=1,2,\dots$).
By Theorem~\ref{thm:trim}, there is a recursive mapping $e\mapsto e'$ with the
property that $\nu^{x}\phi_{e'}(x,y,m,\nu)\le 2^{-m+1}$, and for each $y,m,\nu$ if
$\nu^{x}\phi_{e}(x,y,m,\nu)< 2^{-m+1}$ then $\phi_{e}=\phi_{e'}$.
Let us apply this transformation to the function
$\phi_{e}(x,y,m,\nu)=f_{\nu}(x,y)$.
The result is a lower semicomputable function 
$f'_{\nu}(x,y,m)=\phi_{e'}(x,y,m,\nu)$
with the property that $\nu^{x}f'_{\nu}(x,y,m)\le 2^{-m+1}$, further
$\nu^{x}f_{\nu}(x,y)\le 2^{-m}$ implies $f'_{\nu}(x,y,m)=f_{\nu}(x,y)$.
Now $(x,y,m,\nu)\mapsto 2^{m-1}f'_{\nu}(x,y,m)$ is a uniform test
of $x$ conditional on $y,m$ and hence it is $\lem\t_{\nu}(x \mvert y, m)$.
Substituting $F_{\nu}(y)$ for $m$ the relation
$\nu^{x}f_{\nu}(x,y)\le 2^{-m}$ is satisfied and hence we have
 \begin{align*}
 f_{\nu}(x,y) =
 f'_{\nu}(x,y,F_{\nu}(y)) \lem 2^{-F_{\nu}(y)+1}\t_{\nu}(x\mvert y,F_{\nu}(y)).
 \end{align*}
 \end{proof}

As mentioned above, the
theory generalizes to measures that are not probability measures.
Taking $f_{\mu}(x,y)=1$ in Proposition~\ref{propo:univ-test-gener} gives
the inequality
 \[
   H_{\mu}(x \mvert \flo{\log \mu(X)}) \lea \log\mu(X),
 \]
with a physical meaning when $\mu$ is the phase space measure.
Using~\eqref{eq:H.x.cond.z}, this implies
 \begin{equation}\label{eq:unif.ub}
  H_\mu(x)\lea \log\mu(X) + K(\flo{\log\mu(X)}).
  \end{equation}

% Gacs source line 6451
 \begin{theorem}[General addition]\label{thm:addition-general}
  \uses{def:g4.5.4, thm:g3.1.4}
The following inequality holds:
 \[
   H_{\mu,\nu}(x,y) \eqa
  H_{\mu}(x \mvert \nu)+ H_{\nu}(y \mvert x,\; H_\mu(x \mvert \nu),\; \mu).
 \]
 \end{theorem}

\begin{formalization}
Planned over $\dS$, with $\mu,\nu$ measures of Definition~\ref{our:measure} or $\#$, and with
$H_{\mu}(x \mvert \nu)$ rounded to an integer where it enters as a condition, as
Proposition~\ref{propo:univ-test-gener} requires. For $\mu = \nu = \#$ it is the addition
theorem~\eqref{eq:Haddit}, $K(x,y)\eqa K(x)+K(y \mvert x,K(x))$, formalized in Chapter 3 as
\texttt{addition}; the plan obtains that case through Definition~\ref{our:m-cond}(b).
\end{formalization}
 \begin{proof}
  \proves{thm:addition-general}
  \uses{propo:int.H.of.xy, lem:mon, propo:univ-test-gener}
To prove the inequality $\lea$ define
$G_{\mu,\nu}(x,y,m) =\min_{i\ge m}\;i +  H_{\nu}(y \mvert x, i, \mu)$.
This function is upper semicomputable and increasing
in $m$.
Therefore function
  \[
   G_{\mu,\nu}(x,y)  = G_{\mu,\nu}(x, y, H_{\mu}(x \mvert \nu))
  \]
is also upper semicomputable since it is obtained by substituting an
upper semicomputable function for $m$ in $G_{\mu,\nu}(x,y,m)$.
Lemma~\ref{lem:mon} implies 
  \begin{align*}
     G_{\mu,\nu}(x,y,m) &\eqa m +  H_{\nu}(y \mvert x, m, \mu),
 \\  G_{\mu,\nu}(x,y)   &\eqa H_{\mu}(x \mvert \nu) + 
    H_{\nu}(y \mvert x, H_{\mu}(x \mvert \nu), \mu).
  \end{align*}
   Now, we have
  \begin{align*}
    \nu^{y} 2^{-m - H_{\nu}(y \mvert x, m, \mu)} &\le 2^{-m},
 \\   \nu^{y} 2^{-G_{\mu,\nu}(x,y)} &\lem 2^{-H_{\mu}(x \mvert \mu)}.
  \end{align*}
Integrating over $x$ by $\mu$ gives $\mu^{x}\nu^{y}2^{-G} \lem 1$, implying
$H_{\mu,\nu}(x,y) \lea G_{\mu,\nu}(x,y)$ 
by the minimality property of $H_{\mu,\nu}(x,y)$.
This proves the $\lea$ half of the theorem.

To prove the inequality $\gea$ let $f_{\nu}(x,y,\mu) = 2^{-H_{\mu,\nu}(x,y)}$.
Proposition~\ref{propo:int.H.of.xy} implies that there is a constant $c$ with
$\nu^{y} f_{\nu}(x,y,\mu) \le 2^{-H_{\mu}(x \mvert \nu)+c}$.
Let 
 \begin{align*}
  F_{\nu}(x, \mu) = H_{\mu}(x \mvert \nu).
 \end{align*}
Proposition~\ref{propo:univ-test-gener} gives (substituting $y$ for $x$ and
  $(x,\mu)$ for $y$):
 \begin{align*}
  H_{\mu,\nu}(x,y) = -\log f_{\nu}(x,y,\mu) 
   \gea F_{\nu}(x,\mu)+ H_{\nu}(y \mvert x,F_{\nu}(x,\mu),\mu),
 \end{align*}
which is what needed to be proved.
 \end{proof}

  The function $H_\mu(\og)$ behaves quite differently for different
kinds of measures $\mu$.
Recall the following property of complexity:
 \begin{equation} \label{eq:compl.of.fun}
  K(f(x)\mvert y)\lea K(x\mvert g(y)) \lea K(x) .
 \end{equation}
for any computable functions $f,g$.
This implies
 \[
   K(y)\lea K(x,y).
 \]
  In contrast, if $\mu$ is a probability measure then
 \[
    H_\nu(y) \gea H_{\mu,\nu}(\og,y).
 \]
  This comes from the fact that $2^{-H_\nu(y)}$ is a test for
$\mu\times\nu$.

\subsection{Information}

Mutual information has been defined in Definition~\ref{def:I*}
as $I^{*}(x : y) = K(x) + K(y) - K(x,y)$.
By the Addition theorem, we have
$I^{*}(x:y) \eqa K(y) - K(y \mvert x,\, K(x)) \eqa K(x) - K(x \mvert y,\,K(y))$.   
The two latter expressions show that in some sense, $I^{*}(x:y)$ is the
information held in $x$ about $y$ as well as the information held in $y$
about $x$.  
(The terms $K(x)$, $K(y)$ in the conditions are
logarithmic-sized corrections to this idea.)
Using~\eqref{eq:ol-d}, it is interesting to view
mutual information $I^{*}(x : y)$ as a deficiency of randomness of
the pair $(x,y)$
in terms of the expression $\ol\d_{\mu}$, with respect to $\m \times \m$:
 \[
   I^{*}(x : y) = K(x) + K(y) - K(x,y) = \ol\d_{\m \times \m}(x, y).
 \]
Taking $\m$ as a kind of ``neutral'' probability, even if it is not quite
such, allows us to view $I^{*}(x:y)$ as a ``deficiency of independence''.
Is it also true that $I^{*}(x:y) \eqa \d_{\m \times \m}(x)$?
This would allow us to deduce, as Levin did, ``information conservation''
laws from randomness conservation laws.\footnote{We cannot use the 
test $\ol\t_{\mu}$ for this, since it can be shown easily
that it does not to obey randomness conservation.}

Expression $\d_{\m \times \m}(x)$ must be understood
again in the sense of compactification, as in
Section~\ref{sec:neutral}.
There seem to be two reasonable ways to compactify the space
$\dN\times\dN$: we either compactify it directly, by adding a symbol
$\infty$, or we form the product $\ol\dN \times \ol\dN$.
With either of them, preserving
Theorem~\ref{thm:test-charac-discr}, we would
have to check whether $K(x,y \mvert \m \times \m) \eqa K(x,y)$.
But, knowing the function $\m(x)\times\m(y)$ we 
know the function $x \mapsto \m(x) \eqm \m(x) \times \m(0)$,
hence also the function $(x,y)\mapsto\m(x,y) = \m(\ang{x,y})$.
Using this knowledge, it is possible to develop an argument similar to
the proof of Theorem~\ref{thm:no-upper-semi-neutral}, showing that
$K(x,y \mvert \m \times \m) \eqa K(x,y)$ does not hold.

 \begin{question} 
Is there a neutral measure $M$ with the property that
$I^{*}(x:y) = \d_{M\times M}(x,y)$?
Is this true maybe for all neutral measures $M$?
If not, how far apart are the expressions $\d_{M\times M}(x,y)$ and
$I^{*}(x:y)$ from each other?
 \end{question}


  \section{Randomness and complexity}

We have seen in the discrete case that complexity and randomness are
closely related.
The connection is more delicate technically in the continuous case, but its
exploration led to some nice results.

\subsection{Discrete space}
It is known that for computable \( \mu \), the test \( \d_{\mu}(x) \) can be
expressed in terms of the description complexity of \( x \)
(we will prove these expressions below).
Assume that \( \bX \) is the (discrete) space of all binary strings.
Then we have 
\begin{equation}\label{eq:test-charac-fin-cpt}
   \d_{\mu}(x) = -\log \mu(x) - K(x) + O(K(\mu)).
\end{equation}
The meaning of this equation is the following.  
The expression
\( -\log\mu(x) \) is an upper bound (within \( O(K(\mu)) \)) of the complexity
\( K(x) \), and nonrandomness of \( x \) is measured by
the difference between the complexity and this upper bound.
Assume that \( \bX \) is the space of infinite binary sequences.
Then equation~\eqref{eq:test-charac-fin-cpt} must be replaced with
\begin{equation}\label{eq:test-charac-infin-cpt}
   \d_{\mu}(x) = 
 \sup_{n}\Paren{-\log \mu(x^{\le n}) - K(x^{\le n}) + O(K(\mu))}.
\end{equation}
For noncomputable measures, we cannot replace \( O(K(\mu)) \) in
these relations with anything finite, as shown in
the following example.
Therefore however attractive and simple, 
\( \exp(-\log \mu(x) - K(x)) \) is not a universal uniform test of randomness. 

% Gacs source line 6603
 \begin{proposition}\label{propo:test-charac-counterexample}
  \uses{def:g4.1.5}
There is a measure \( \mu \) over the discrete space \( \bX \) of binary strings
such that for each \( n \), there is an \( x \) with
\( \d_{\mu}(x) = n - K(n) \) and \( -\log \mu(x) - K(x) \lea 0 \).
 \end{proposition}

\begin{formalization}
Planned. The measure $\mu$ lives on the pairs $(0^{n},y_{n})$ and is not computable; a $y_{n}$
with $K(0^{n},y_{n})>n$ exists by counting. The test $t_{\mu}$ is lower semicomputable in
$(x,\mu)$, since its denominator is a finite sum of values of $\mu$.
\end{formalization}
 \begin{proof}
  \proves{propo:test-charac-counterexample}
  \uses{our:lsc}
Let us treat the domain of our measure \( \mu \) as a set of pairs \( (x,y) \).
Let \( x_{n} = 0^{n} \), for \( n=1,2,\dotsc \).
For each \( n \), let \( y_{n} \) be some binary string of length \( n \)
with the property \( K(x_{n}, y_{n}) > n \).
Let \( \mu(x_{n},y_{n})=2^{-n} \), and 0 elsewhere.
Then \( - \log \mu(x_{n},y_{n}) - K(x_{n},y_{n}) \le n - n = 0 \).
On the other hand, let \( t_{\mu}(x,y) \) be the test nonzero only on pairs of strings
\( (x,y) \) of the form \(( x_{n},y) \):
 \[
 t_{\mu}(x_{n}, y) = \frac{\m(n)}{\sum_{z \in \cB^{n}} \mu(x_{n}, z)}.
 \]
The form of the definition ensures semicomputability and we also have
 \[
 \sum_{x,y} \mu(x,y) t_{\mu}(x,y) \le \sum_{n} \m(n) < 1,
 \]
therefore \( t_{\mu} \) is indeed a test.
Hence \( \t_{\mu}(x,y) \gem t_{\mu}(x,y) \).
Taking logarithms, \( \d_{\mu}(x_{n}, y_{n}) \gea n - K(n) \).
 \end{proof}

 The same example shows that the test defined as
\( \exp(-\log\mu(x) - K(x)) \) over discrete sets,
does not satisfy the randomness conservation property.

% Gacs source line 6633
 \begin{proposition}
  \label{prop:g4.6.2}
  \uses{propo:test-charac-counterexample}
The test defined as \( f_{\mu}(x) = \exp(-\log\mu(x) - K(x)) \) 
over discrete spaces \( \bX \) does
not obey the conservation of randomness.
 \end{proposition}
 \begin{proof}
  \proves{prop:g4.6.2}
  \uses{propo:test-charac-counterexample}
Let us use the example of Proposition~\ref{propo:test-charac-counterexample}.
Consider the function \( \pi : (x,y) \mapsto x \).
The image of the measure \( \mu \) under the projection is 
\( (\pi\mu)(x) = \sum_{y} \mu(x,y) \).
Thus, \( (\pi\mu)(x_{n}) = \mu(x_{n},y_{n}) = 2^{-n} \).
Then we have seen that \( \log f_{\mu}(x_{n},y_{n}) \le 0 \).
On the other hand,
 \[ 
 \log f_{\pi\mu} (\pi(x_{n},y_{n})) = -\log(\pi\mu)(x_{n}) - K(x_{n})
 \eqa n - K(n). 
 \]
Thus, the projection \( \pi \) takes a random pair \( (x_{n},y_{n}) \) into
an object \( x_{n} \) that is very nonrandom (when randomness is measured using
the test \( f_{\mu} \)).
 \end{proof}
In the example, we have the abnormal situation that a pair is random but
one of its elements is nonrandom.
Therefore even if we would not insist on universality, the test 
\( \exp(-\log\mu(x) - K(x)) \) is unsatisfactory.

Looking into the reasons of the nonconservation in the example,
we will notice that it could only have happened because the 
test \( f_{\mu} \) is too special.
The fact that \( -\log (\pi\mu)(x_{n}) - K(x_{n}) \) is large should show that 
the pair \( (x_{n},y_{n}) \) can be enclosed into the ``simple'' set
\( \{x_{n}\} \times \bY \) of small probability; unfortunately, 
this observation does not reflect on 
\( -\log\mu(x,y) - K(x,y) \) (it does for computable \( \mu \)).
 
It is a natural idea to modify equation~\eqref{eq:test-charac-fin-cpt}
in such a way that the complexity \( K(x) \) is replaced with \( K(x \mvert \mu) \).
However, this expression must be understood properly.
We need to use the definition of \( K(x) \) as \( -\log\m(x) \) directly, and not 
as prefix complexity.

Let us mention the following easy fact:

% Gacs source line 6678
 \begin{proposition}\label{propo:H/mu-computable}
  \uses{our:m-cond, def:m, thm:univsem, our:lsc-props}
If \( \mu \) is a computable measure then \( K(x \mvert \mu) \eqa K(x) \).
The constant in \( \eqa \) depends on the description complexity of \( \mu \).
 \end{proposition}

\begin{formalization}
Planned, with the constant depending on $\mu$. At a computable $\mu$, $\m(x \mvert \mu)$ is lower
semicomputable in $x$ (Theorem~\ref{our:lsc-props}(d)), hence $\lem\m(x)$; and $\m(x)$, constant
in $\mu$, is lower semicomputable in $(x,\mu)$, hence $\lem\m(x \mvert \mu)$.
\end{formalization}

% Gacs source line 6683
 \begin{theorem}\label{thm:test-charac-discr}
  \uses{def:g4.1.5, our:m-cond, thm:univ-integrable-test}
If \( \bX \) is the discrete space \( \Sg^{*} \) then we have
 \begin{equation}\label{eq:test-charac-discr}
   \d_{\mu}(x) \eqa -\log\mu(x) - K(x \mvert \mu).
 \end{equation}
 \end{theorem}

\begin{formalization}
Planned over $\dS$, as $\t_{\mu}(x)\eqm\m(x \mvert \mu)/\mu(x)$ in $\clint{0}{\infty}$. For a
computable $\mu$, by Proposition~\ref{propo:H/mu-computable} and the coding theorem, this is
Theorem~2.2.1 (\texttt{canonicalIntegrableTest\_isUniversal},
\texttt{universalSemimeasure\_div\_isUniversal}), and the proof is the same: $\mu(x)\t_{\mu}(x)$
and $\m(x \mvert \mu)/\mu(x)$ are lower semicomputable in $(x,\mu)$ because $\mu(x)$ is
computable from a name of $\mu$.
\end{formalization}
Note that in terms of the algorithmic entropy notation introduced 
in~\eqref{eq:alg-ent}, this theorem can be expressed as
 \[
   H_{\mu}(x) \eqa K(x \mvert \mu) + \log\mu(x).
 \]
 \begin{proof}
  \proves{thm:test-charac-discr}
  \uses{our:lsc-props}
In exponential notation, equation~\eqref{eq:test-charac-discr} can be
written as \( \t_{\mu}(x) \eqm \m(x \mvert \mu)/\mu(x) \).
Let us prove \( \gem \) first.
We will show that the right-hand side of this inequality is a test, and
hence \( \lem \t_{\mu}(x) \).
However, 
the right-hand side is clearly lower semicomputable in \( (x, \mu) \) and
when we ``integrate'' it (multiply it by \( \mu(x) \) and sum it), its sum is
\( \le 1 \); thus, it is a test.

Let us prove \( \lem \) now.
The expression \( \t_{\mu}(x)\mu(x) \) is clearly lower semicomputable in
\( (x,\mu) \), and its sum is \( \le 1 \).
Hence, it is \( \lea \m(x \mvert \mu) \).
 \end{proof}

% Gacs source line 6711
 \begin{remark}\label{rem:oracle-warning}
It important not to consider relative computation with respect to \( \mu \) 
as \emph{oracle computation} in the ordinary sense.
Theorem~\ref{thm:neutral-meas} below will show the existence of a measure with
respect to which every element is random.
If randomness is defined using \( \mu \) as an oracle then we can always find
elements nonrandom with respect to \( \mu \).

For similar reasons, the proof of the Coding Theorem does not transfer to the
function \( K(x\mvert \mu) \) 
since an interpreter function should have the property of \df{intensionality},
depending only on \( \mu \) and not on the sequence representing it.
(It does transfer without problem to an oracle version of \( K^{\mu}(x) \).)
The Coding Theorem still may hold, at least in some cases: this is currently not
known.
Until we know this, we cannot interpret \( K(x\mvert \mu) \) as description complexity
in terms of interpreters and codes.

(Thanks to Alexander Shen for this observation: this remark corrects
an error in the paper~\cite{GacsUnif05}.)
 \end{remark}


\begin{formalization}
This is why Definition~\ref{our:m-cond} defines $\m(x \mvert \mu)$ as the maximal semimeasure
that is lower semicomputable in $(x,\mu)$, independently of names, and not through Chapter 3's
oracle machines (\texttt{KPrel}), whose value may depend on the name. Theorem~\ref{our:oracle-bound}
records the inequality between the two that does hold.
\end{formalization}

% Ours: the one inequality between oracle complexity and K(x | mu)
\begin{ourtheorem}[Oracle complexity and \( K(x \mvert \mu) \)]\label{our:oracle-bound}
  \uses{our:m-cond, our:measure, our:lsc-props}
Let \( \bX \) be \( \dS \) or \( \dB^{\dN} \).
There is a constant \( c \) such that for every measure \( \mu\in\cM(\bX) \), every name \( A \)
of \( \mu \) and every string \( x \),
 \[
   K^{A}(x) \le K(x \mvert \mu) + c,
 \]
where \( K^{A}(x) \) is the prefix complexity of \( x \) relative to the oracle \( A \).
\end{ourtheorem}

\begin{formalization}
Planned. $K^{A}(x)$ is Chapter 3's \texttt{KPrel V A x} for an optimal oracle machine
\texttt{V} (\texttt{IsOptimalOracle}, \texttt{exists\_isOptimalOracle}). Proof: by
Theorem~\ref{our:lsc-props}(a), finite prefixes of $A$ yield lower bounds for
$\m(x \mvert \mu)$ that are stable under extension of the prefix. Issued as requests to the
relativized Kraft--Chaitin coder \texttt{exists\_oracleMachine\_of\_requests}, and cut off where
their Kraft sum would exceed 1 (which does not happen along a name), they give the bound. The
reverse inequality cannot be expected (Remark~\ref{rem:oracle-warning}). Nothing else in the
chapter uses this theorem.
\end{formalization}


\subsection{Non-discrete spaces}
For non-discrete spaces, unfortunately, we can only provide a less
intuitive expression.

% Gacs source line 6738
 \begin{proposition}\label{thm:test-charac}
  \notready
let \( \bX=(X, d, D, \ag) \) be a complete
computable metric space, and let \( \cE \) be the enumerated set of bounded
Lipschitz functions introduced in~(A.1.3), 
but for the space \( \bM(\bX) \times \bX \).
The uniform test of randomness \( \t_{\mu}(x) \), can be expressed as
 \begin{equation}\label{eq:test-charac}
  \t_{\mu}(x) \eqm 
  \sum_{f \in \cE}f(\mu,x)\frac{\m(f \mvert \mu)}{\mu^{y} f(\mu, y)}.
 \end{equation}
 \end{proposition}

\begin{formalization}
Not planned: the sum ranges over the bounded Lipschitz functions $\cE$ of a general computable
metric space (A.1.3).
\end{formalization}
 \begin{proof}
  \proves{thm:test-charac}
  \notready
For \( \gem \),
we will show that the right-hand side of the inequality is a test, and
hence \( \lem \t_{\mu}(x) \).
For simplicity, we skip the notation about the enumeration of \( \cE \) and
treat each element \( f \) as its own name.
Each term of the sum is clearly lower semicomputable in
\( (f, x, \mu) \), hence the sum is lower semicomputable in \( (x, \mu) \).
It remains to show that the \( \mu \)-integral of the sum is \( \le 1 \).
But, the \( \mu \)-integral of the generic term is \( \le \m(f \mvert \mu) \), and
the sum of these terms is \( \le 1 \) by the definition of the function
\( \m(\cdot \mvert \cdot) \).
Thus, the sum is a test.

For \( \lem \), note that \( (\mu,x) \mapsto \t_{\mu}(x) \), 
as a lower semicomputable function, is the supremum of functions in \( \cE \).
Denoting their differences by \( f_{i}(\mu,x) \),
we have \( \t_{\mu}(x) = \sum_{i} f_{i}(\mu,x) \).
The test property implies \( \sum_{i} \mu^{x} f_{i}(\mu,x) \le 1 \).
Since the function \( (\mu,i) \mapsto \mu^{x} f_{i}(\mu,x) \) is lower
semicomputable, this implies \( \mu^{x} f_{i}(\mu,x) \lem \m(i \mvert \mu) \), and
hence 
 \[
 f_{i}(\mu,x) \lem f_{i}(\mu,x) \frac{\m(i \mvert \mu)}{\mu^{x} f_{i}(\mu,x)}.
 \]
It is easy to see that for each \( f\in\cE \) we have
 \[
   \sum_{i : f_{i} = f} \m(i \mvert \mu) \le \mu(f \mvert \mu),
 \]
which leads to~\eqref{eq:test-charac}.
 \end{proof}

 \begin{remark}\label{rem:test-charac-lb}
If we only want the \( \gem \) part of the result, then \( \cE \) can be
replaced with any enumerated computable sequence of bounded computable
functions.
 \end{remark}


  \subsection{Infinite sequences}

In case of the space of infinite sequences and a computable measure,
Theorem~\ref{thm:defic-charac-cpt-seqs} gives a characterization of
randomness in terms of complexity.
This theorem does not seem to transfer to a more general situation, but
under some conditions, at least parts of it can be extended.

For arbitrary measures and spaces, we can say a little less:

% Gacs source line 6798
 \begin{proposition}\label{propo:test-charac-seq-lb}
  \uses{def:g4.1.5, our:m-cond, thm:defic-charac-cpt-seqs}
For all measures \( \mu\in\cM_{R}(X) \), 
for the deficiency of randomness \( \d_{\mu}(x) \), we have
 \begin{equation}\label{eq:defic-ineq-seq}
 \d_{\mu}(x) \gea \sup_{n}\Paren{-\log \ol\mu(x^{\le n}) - K(x^{\le n} \mvert \mu)}.
 \end{equation}
 \end{proposition}

\begin{formalization}
Planned over $\dB^{\dN}$, where the cells $\Gg_{s}$ are the cylinders, $\ol\mu = \mu$ for a
probability measure, and $\cM_{1}(X)$ is Definition~\ref{our:measure}. It generalizes the
$\gea$ half of Theorem~2.3.4 (\texttt{canonicalSeqTest\_isIntegrable}) to noncomputable $\mu$,
with $K(\cdot \mvert \mu)$ in place of $K$.
\end{formalization}
 \begin{proof}
  \proves{propo:test-charac-seq-lb}
  \uses{our:lsc-props}
  Consider the function
 \[
 f_{\mu}(x) = \sum_{s} 1_{\Gg_{s}}(x) \frac{\m(s \mvert \mu)}{\ol\mu(\Gg_{s})}
  = \sum_{n} \frac{\m(x^{\le n} \mvert \mu)}{\ol\mu(x^{\le n})}
  \ge \sup_{n} \frac{\m(x^{\le n} \mvert \mu)}{\ol\mu(x^{\le n})}.
 \]
The function \( (\mu,x) \mapsto f_{\mu}(x) \) 
is clearly lower semicomputable and satisfies
\( \mu^{x} f_{\mu}(x) \le 1 \), and hence 
 \[
  \d_{\mu}(x) \gea \log f(x) \gea 
   \sup_{n}\Paren{-\log\ol\mu(x^{\le n}) - K(x^{\le n} \mvert \mu)}.
 \]
 \end{proof}

% Gacs source line 6821
\begin{definition}
  \label{def:g4.6.8}
  \uses{our:measure}
Let \( \cM_{R}(X) \) be the set of measures \( \mu \) with \( \mu(X)=R \).  
\end{definition}

\begin{formalization}
In the plan $R = 1$, and $\cM_{1}(X)$ is Definition~\ref{our:measure}.
\end{formalization}

We will be able to prove the \( \gea \) part of the statement of
Theorem~\ref{thm:defic-charac-cpt-seqs} in a more general
space, and without assuming computability.
Assume that a separating sequence \( b_{1},b_{2},\dots \) is given as defined
in Subsection~\ref{subsec:cells}, along with the set \( X^{0} \).
For each \( x \in X^{0} \), the binary sequence \( x_{1},x_{2},\dots \) 
has been defined.
Let
 \begin{align*}
  \ol\mu(\Gg_{s}) &= R  - \sum\setof{\mu(\Gg_{s'}) : \len{s}=\len{s'},\;s'\ne s}.
 \end{align*}
Then \( (s,\mu)\mapsto \mu(\Gg_{s}) \) is lower semicomputable, and 
\( (s,\mu)\mapsto \ol\mu(\Gg_{s}) \) is upper semicomputable.
And, every time that the functions \( b_{i}(x) \) form a partition
with \( \mu \)-continuity, we have \( \ol\mu(\Gg_{s})=\mu(\Gg_{s}) \) for all \( s \).

% Gacs source line 6841
 \begin{theorem}\label{thm:defic-charac-compact}
  \uses{def:g4.6.8, def:g4.1.5, thm:defic-charac-cpt-seqs}
Suppose that the space \( X \) is effectively compact.
Then for all computable measures \( \mu\in\cM_{R}^{0}(X) \), 
for the deficiency of randomness \( \d_{\mu}(x) \), 
the characterization~\eqref{eq:defic-charac-seq} holds.
 \end{theorem}

\begin{formalization}
Planned only for $X = \dB^{\dN}$ with the cylinders as cells ($X^{0} = X$, $\cM^{0}_{1} = \cM$);
general effectively compact spaces need Section~\ref{subsec:cells} and are out of scope. There it
is Theorem~2.3.4, formalized in Chapter 2 (\texttt{universalSeqTest\_equiv}) for Chapter 2's
universal test. The plan adds that for computable $\mu$ the uniform test $\t_{\mu}$ agrees with
that test up to a constant factor: each is lower semicomputable in $\xi$ and a $\mu$-test, so
each dominates the other (Corollary~\ref{coroll:trim} in one direction).
\end{formalization}
 \begin{proof}
  \proves{thm:defic-charac-compact}
  \uses{propo:test-charac-seq-lb, coroll:trim, thm:univ-integrable-test-seq}
The proof of part \( \gea \) of the inequality follows directly from
Proposition~\ref{propo:test-charac-seq-lb}, just as in the proof of 
Theorem~\ref{thm:defic-charac-cpt-seqs}.

The proof of \( \lea \) is also similar to the proof of that theorem.
The only part that needs to be reproved is the statement that
for every lower semicomputable function \( f \) over \( X \),
there are computable sequences \( y_{i}\in\dN^{*} \) and
\( q_{i}\in\dQ \) with \( f(x) = \sup_{i} q_{i} 1_{y_{i}}(x) \).
This follows now, since according to
Proposition~\ref{propo:compact-cell-basis}, the cells \( \Gg_{y} \) form a basis of
the space \( X \).
 \end{proof}


\subsection{Bernoulli tests}

In this part, we will give characterize Bernoulli sequences in terms of their
complexity growth.

Recall Definition~\ref{def:combinat-Bernoulli}:
A function \( f:\dB^{*}\to\dR \) is a combinatorial Bernoulli test if
 \begin{alphenum}
  \item It is lower semicomputable.
  \item It is monotonic with respect to the prefix relation.
  \item For all \( 0\le k\le n \) we have \( \sum_{x\in\dB(n,k)} f(x) \le \binom{n}{k} \).
 \end{alphenum}
According to Theorem~\ref{thm:combinat-Bernoulli}, a 
sequence \( \xi \) is nonrandom with respect to all Bernoulli measures if and
only if \( \sup_{n}b(\xi^{\le n})=\infty \), where \( b(x) \) is a combinatorial
Bernoulli test.

We need some definitions.

% Gacs source line 6882
\begin{definition}
  \label{def:g4.6.9}
  \uses{def:semicomputability}
For a finite or infinite sequence \( x \) let \( S_{n}(x)=\sum_{i=1}^{n}x(i) \).

For \( 0\le p\le 1 \) and integers \( 0\le k\le n \), 
denote \( B_{p}(n,k)=\binom{n}{k}p^{k}(1-p)^{n-k} \).

An upper semicomputable function \( D:\dN^{2}\to\dN \), defined for \( n\ge 1 \),
\( 0\le k\le n \) will be called a \df{gap function} if
 \begin{align}\label{eq:dg-bd}
  \sum_{n\ge 1}\sum_{k=0}^{n}B_{p}(n,k)2^{-D(n,k)} \le 1
 \end{align}
holds for all \( 0\le p\le 1 \).
A gap function \( D(n,k) \) is \df{optimal} if for every other gap 
function \( D'(n,k) \) there is a \( c_{D'} \) with \( D(n,k)\le D'(n,k)+c_{D'} \).
\end{definition}

\begin{formalization}
Planned. Only numbers are involved: $D$ is upper semicomputable on $\dN^{2}$ in Chapter 1's sense,
and~\eqref{eq:dg-bd} is required for every real $p\in\clint{0}{1}$.
\end{formalization}

% Gacs source line 6898
 \begin{proposition}
  \label{prop:g4.6.10}
  \uses{def:g4.6.9}
There is an optimal gap function \( D(n,k)\lea K(n) \).
 \end{proposition}

\begin{formalization}
Planned. The trimming needs to recognize $\sup_{p}\sum B_{p}(n,k)2^{-D(n,k)}<1$ for finite sums;
the supremum of a polynomial in $p$ over $\clint{0}{1}$ is computable.
\end{formalization}
 \begin{proof}
  \proves{prop:g4.6.10}
  \uses{thm:trim, l:kraft}
The existence is proved using the technique of Theorem~\ref{thm:trim}.
For the inequality it is sufficient to note that \( K(n) \) is a gap function.
Indeed, we have
 \begin{align*}
  \sum_{n\ge 1}\sum_{k=0}^{n}B_{p}(n,k)2^{-K(n)} =\sum_{n\ge 1}2^{-K(n)}\le 1.
 \end{align*}
 \end{proof}

% Gacs source line 6911
\begin{definition}
  \label{def:g4.6.11}
  \uses{prop:g4.6.10}
Let us fix an optimal gap function and denote it by \( \Dg(n,k) \).
\end{definition}

Now we can state the test characterization theorem for Bernoulli tests.

% Gacs source line 6917
 \begin{theorem}
  \label{thm:g4.6.3}
  \uses{def:g4.6.11, def:g4.2.14, our:m-cond}
Denoting by \( b(\xi) \) the universal class test for the Bernoulli sequences, we
have \( b(\xi) \eqm \ol b(\xi) \),
where
  \begin{align*}
  \log \ol b(\xi) = \sup_{n}\log\binom{n}{k}-K(\xi^{\le n}\mvert n,k,\Dg(n,k))-\Dg(n,k),
 \end{align*}
with \( k=S_{n}(\xi) \).
 \end{theorem}

\begin{formalization}
Planned. Everything here is on strings and numbers: the conditions $(n,k,\Dg(n,k))$ are discrete,
and $\m(y \mvert n,k,\Dg(n,k))$ is Chapter 1's conditional universal semimeasure
(Definition~\ref{our:m-cond}(b) with a string condition). The decomposition
$b'(\xi) = \sup_{i}2^{k_{i}}1_{y_{i}}(\xi)$ is the one Chapter 2 made decidable for
Theorem~2.3.4 (first hits of dyadic levels, \texttt{LevinSchnorr.lean}).
\end{formalization}
 \begin{proof}
  \proves{thm:g4.6.3}
  \uses{propo:univ-test-gener, thm:combinat-Bernoulli}
Let \( \dg(n,k) = 2^{-\Dg(n,k)} \).
 \begin{claim}\label{claim:gap}
Consider lower semicomputable functions \( \gm:\dB^{*}\to\dR_{+} \) such that
for all \( n,k \) we have
 \begin{align*}
   \sum_{y\in\dB(n,k)}\gm(y) \le 2^{-\Dg(n,k)}.
 \end{align*}
Among these functions there is one
that is optimal (maximal to within a multiplicative constant).
Calling it \( \dg(y) \) we have
 \begin{align*}
  \dg(y) \eqm \dg(n,k)\cdot\m(y\mvert n,k,\Dg(n,k)).
 \end{align*}
 \end{claim}
Thus, the right-hand side is equal, within a multiplicative constant,
to a lower semicomputable function of \( y \).
 \begin{proof}
  \proves{claim:gap}
This follows immediately from Proposition~\ref{propo:univ-test-gener}.
with \( \nu=\# \) (the counting measure) over the sets \( \dB(n,k) \).
 \end{proof}
Let us show \( \ol b(\xi)\lem b(\xi) \).
We have with \( k=S_{n}(\xi) \) in the first line:
 \begin{align*}
 \ol b(\xi) &=\sup_{n\ge 1}\binom{n}{k}\m(\xi^{\le n}\mvert n,k,\Dg(n,k))\dg(n,k)
\\ &= \sup_{n\ge 1}\sum_{k=0}^{n}\binom{n}{k}\dg(n,k)
     \sum_{y\in\dB(n,k)}1_{y}(\xi)\m(y\mvert n,k,\Dg(n,k))
\\ &\le \sum_{n\ge 1}\sum_{k=0}^{n}\binom{n}{k}
     \sum_{y\in\dB(n,k)}1_{y}(\xi)\dg(n,k)\m(y\mvert n,k,\Dg(n,k))
\\ &\eqm \sum_{n\ge 1}\sum_{k=0}^{n}\binom{n}{k}
     \sum_{y\in\dB(n,k)}1_{y}(\xi)\dg(y),
 \end{align*}
using the notation of Claim~\ref{claim:gap} above.
Let \( t(\xi) \) denote the right-hand side here, which is thus
a lower semicomputable function.
We have for all \( p \):
 \begin{align*}
   B_{p}^{\xi}t(\xi) &\eqm \sum_{n\ge 1}\sum_{k=0}^{n}B_{p}(n,k)\dg(n,k)
 \sum_{y\in\dB(n,k)}\m(y\mvert n,k,\Dg(n,k)) \le 1,
 \end{align*}
so \( t(\xi)\gem \ol b(\xi) \) is a Bernoulli test, showing \( \ol b(\xi)\lem b(\xi) \).

To show \( b(\xi)\lem \ol b(\xi) \) we will
follow the method of the proof of
Theorem~\ref{thm:defic-charac-cpt-seqs}.
Replace \( b(\xi) \) with a rougher version:
 \[
 b'(\xi) = \frac{1}{2}\max \setof{2^{n} : 2^{n} < b(\xi)},
 \]
then we have \( 2 b'<b \).
There are computable sequences \( y_{i}\in\dB^{*} \) and \( k_{i} \in \dN \) with
\( b'(\xi) = \sup_{i}\; 2^{k_{i}} 1_{y_{i}}(\xi) \)
with the property that if \( i<j \) and \( 1_{y_{i}}(\xi)=1_{y_{j}}(\xi)=1 \) then
\( k_{i}<k_{j} \).
As in the imitated proof, we have
\( 2b'(\xi) \ge \sum_{i} 2^{k_{i}} 1_{y_{i}}(\xi) \).
The function \( \gm(y)=\sum_{y_{i}=y}2^{k_{i}} \) is lower semicomputable.
We have
\begin{align}
 b\ge 2b'(\xi)\ge\sum_{i} 2^{k_{i}} 1_{y_{i}}(\xi)  = \sum_{y\in\dN^{*}}1_{y}(\xi)\gm(y)
 = \sum_{n}\sum_{k=0}^{n}\sum_{y\in\dB(n,k)}1_{y}(\xi)\gm(y).
\label{eq:gamma-expr}
 \end{align}
By Theorem~\ref{thm:combinat-Bernoulli} we can assume
\( \sum_{y\in\dB(n,k)}\gm(y)\le \binom{n}{k} \).
Let
 \begin{alignat*}{3}
           &\dg'(y)    &&= \gm(y)\binom{n}{k}^{-1} \le 1,
\\      &\dg'(n,k) &&=  \sum_{y\in\dB(n,k)}\dg'(y).
 \end{alignat*}
Since \( 1\ge B_{p}b \ge B_{p}(2 b') \) for all \( p \), we have 
 \begin{align*}
      1 &\ge \sum_{n}\sum_{k=0}^{n}\sum_{y\in\dB(n,k)}\gm(y)B_{p}^{\xi}1_{y}(\xi)
            = \sum_{n}\sum_{k=0}^{n}p^{k}(1-p)^{n-k}\sum_{y\in\dB(n,k)}\gm(y)
\\     &= \sum_{n}\sum_{k=0}^{n}B_{p}(n,k)\sum_{y\in\dB(n,k)}\dg'(y)
          = \sum_{n}\sum_{k=0}^{n}B_{p}(n,k)\dg'(n,k).
 \end{align*}
Thus \( \dg'(n,k) \) is a gap function, hence \( \dg'(n,k) \lem 2^{-\Dg(n,k)} \),
and by Claim~\ref{claim:gap} we have 
 \begin{align*}
 \gm(y)\binom{n}{k}^{-1}&=\dg'(y)\lem \dg(y)
   \eqm\dg(n,k)\cdot\m(y\mvert n,k,\Dg(n,k)).
 \end{align*}
Substituting back into~\eqref{eq:gamma-expr} finishes the proof of
\( b(\xi)\lem \ol b(\xi) \).
 \end{proof}


\section{Cells}\label{subsec:cells}

This section allows to transfer some of the results
on sequence spaces to more general spaces, by encoding the elements into
sequences.
The reader who is only interested in sequences can skip this section.

\begin{formalization}
Not planned. Cells encode the points of a general space as binary sequences. In the plan the
space already is $\dB^{\dN}$, whose canonical cells are the cylinders, with $X^{0} = X$
(Examples~\ref{example:cells}, part 1), so nothing in this section is needed. Its nodes are
marked not ready.
\end{formalization}

\subsection{Partitions}

The coordinates of the sequences into which we want to encode our elements will
be obtained via certain partitions.

Recall from Definition~A.2.33 that 
a measurable set \( A \) is said to be a \df{\( \mu \)-continuity set} if 
\( \mu(\partial A)=0 \) where \( \partial A \) is the boundary of \( A \).

% Gacs source line 7030
 \begin{definition}[Continuity partition]
  \label{def:g4.7.1}
  \notready
Let \( f:X\to\dR \) be a bounded computable function, and let
\( \ag_{1}<\dots<\ag_{k} \) be rational numbers,
and let \( \mu \) be a computable measure with the property
that \( \mu f^{-1}(\ag_{j})=0 \) for all \( j \).


In this case, we will say that \( \ag_{j} \) are \( \mu \)-\df{continuity points} of \( f \).
Let \( \ag_{0}=-\infty \), \( \ag_{k+1}=\infty \), and for \( j=0,\dots,k \), let
Let \( U_{j} = f^{-1}((j,j+1)) \).
The sequence of disjoint 
computably enumerable open sets \( (U_{0},\dots,U_{k}) \) will be called the
\df{partition generated by} \( f,\ag_{1},\dots,\ag_{k} \).

If we have several partitions \( (U_{i0},\dots,U_{i,k}) \), 
generated by different functions \( f_{i} \) (\( i=1,\dots,m \))
and cutoff sequences \( (\ag_{ij}: j=1,\dots,k_{i}) \)  made up of
\( \mu \)-continuity points of \( f_{i} \)
then we can form a new partition generated by all possible intersections
 \[
  V_{j_{1},\dots,j_{n}} = U_{1,j_{1}}\cap \dots \cap U_{m,j_{m}}.
 \]
A partition of this kind will be called a \df{continuity partition}.
The sets \( V_{j_{1},\dots,j_{n}} \) will be called the \df{cells} of this
partition. 
 \end{definition}

The following is worth noting.

% Gacs source line 7059
 \begin{proposition}\label{propo:reg-partit-meas-cptable}
  \notready
In a partition as given above, 
the values \( \mu V_{j_{1},\dots,j_{n}} \) are computable
from the names of the functions \( f_{i} \) and the cutoff points \( \ag_{ij} \).
 \end{proposition}
 \begin{proof}
  \proves{propo:reg-partit-meas-cptable}
  \notready
Straightforward.
 \end{proof}

Let us proceed to defining cells.

% Gacs source line 7070
 \begin{definition}[Separating sequence]
  \label{def:g4.7.3}
  \notready
Assume that a computable sequence of functions 
\( b_{1}(x),b_{2}(x),\dots \) over \( X \) is given,
with the property that for every pair 
\( x_{1},x_{2}\in X \) with \( x_{1}\ne x_{2} \), there is a \( j \)
with \( b_{j}(x_{1})\cdot b_{j}(x_{2}) < 0 \).
Such a sequence will be called a \df{separating sequence}.
Let us give the correspondence between the set \( \dB^{\dN} \)
of infinite binary sequences and elements of the set
 \[
  X^{0} = \setof{x\in X: b_{j}(x)\ne 0,\;j=1,2,\dots}.
 \]
For a binary string 
\( s_{1}\dotsm s_{n} = s\in\dB^{*} \), let
 \[
   \Gg_{s}
 \]
be the set of elements of \( X \) with the property that for
\( j=1,\dots,n \), if \( s_{j}=0 \) then \( b_{j}(\og) < 0 \), otherwise
\( b_{j}(\og)>0 \).

The separating sequence will be called \( \mu \)-\df{continuity} sequence
if \( \mu(X^{0})=0 \).
 \end{definition}

This correspondence has the following properties.
 \begin{alphenum}
  \item \( \Gg_{\Lg}=X \).
  \item For each \( s\in \dB \), the sets \( \Gg_{s0} \) and 
\( \Gg_{s1} \) are disjoint and their union is contained in \( \Gg_{s} \).
  \item For \( x\in X^0 \), we have \( \{x\} = \bigcap_{x\in\Gg_{s}} \Gg_{s} \).
 \end{alphenum}

% Gacs source line 7103
 \begin{definition}[Cells]
  \label{def:g4.7.4}
  \notready
If string \( s \) has length \( n \) then \( \Gg_{s} \) will be called a \df{canonical
\( n \)-cell}, or simply canonical cell, or \( n \)-cell.
From now on, whenever \( \Gg \) denotes a subset of \( X \), it means a
canonical cell.
We will also use the notation
 \[
   \len{\Gg_{s}}=\len{s}.
 \]
 \end{definition}

The three properties above say that if we restrict ourselves to the set
\( X^{0} \) then the canonical cells behave somewhat like binary subintervals:
they divide \( X^{0} \) in half, then each half again in half, etc.  
Moreover, around each point, these canonical cells become ``arbitrarily
small'', in some sense (though, they may not be a basis of neighborhoods).
It is easy to see that if \( \Gg_{s_1},\Gg_{s_2} \) are two canonical
cells then they either are disjoint or one of them contains the other.
If \( \Gg_{s_1}\sbs\Gg_{s_2} \) then \( s_2 \) is a prefix of \( s_1 \).
If, for a moment, we write \( \Gg^0_s=\Gg_s\cap X^0 \) then we have the
disjoint union \( \Gg^0_s=\Gg^0_{s0}\cup\Gg^0_{s1} \).

Let us use the following notation.

% Gacs source line 7127
 \begin{definition}
  \label{def:g4.7.5}
  \notready
For an \( n \)-element binary string \( s \), for \( x \in \Gg_{s} \), we will write 
 \[
           s = x^{\le n} =  x_1\dotsm x_n,
\quad \mu(s) = \mu(\Gg_{s}).
 \]
The \( 2^n \) cells (some of them possibly empty)
of the form \( \Gg_s \) for \( \len{s}=n \) form a partition
 \[
  \cP_n
 \]
  of \( X^0 \).
 \end{definition}

  Thus, for elements of \( X^0 \), we can talk about the \( n \)-th bit \( x_n \)
of the description of \( x \): it is uniquely determined.

 \begin{examples}\label{example:cells}\
 \begin{enumerate}
  \item If \( \bX \) is the set of infinite binary sequences with its usual
topology, the functions \( b_{n}(x) = x_{n}-1/2 \) generate the usual
cells, and \( \bX^{0}=\bX \).
  \item If \( \bX \) is the interval \( \clint{0}{1} \), let 
\( b_{n}(x) = -\sin(2^{n}\pi x) \).
Then cells are open 
intervals of the form \( \opint{k\cdot 2^{-n}}{(k+1)\cdot 2^{n}} \),
the correspondence between infinite binary strings
and elements of \( X^0 \) is just the usual representation of \( x \) as the
binary decimal string \( 0.x_{1}x_{2}\dots \).
 \end{enumerate}
 \end{examples}

When we fix canonical cells,  
we will generally assume that the partition chosen is also
``natural''.
The bits \( x_1,x_2,\ldots \) could contain information about the
point \( x \) in decreasing order of importance from a macroscopic
point of view. 
For example, for a container of gas, the first few bits may
describe, to a reasonable degree of precision, the amount of gas in
the left half of the container, the next few bits may describe the
amounts in each quarter, the next few bits may describe the
temperature in each half, the next few bits may describe again the
amount of gas in each half, but now to more precision, etc.
From now on, whenever \( \Gg \) denotes a subset of \( X \), it means a
canonical cell.
From now on, for elements of \( X^0 \), we can talk about the \( n \)-th
bit \( x_n \) of the description of \( x \): it is uniquely determined.

The following observation will prove useful.

% Gacs source line 7178
 \begin{proposition}\label{propo:compact-cell-basis}
  \notready
Suppose that the space \( \bX \) is effectively compact\footnote{It was noted by
  Hoyrup and Rojas that the qualification ``effectively'' is necessary here.}
 and we have a separating sequence \( b_{i}(x) \) as given above.
Then the cells \( \Gg_{s} \) form a basis of the space \( \bX \).
 \end{proposition}
 \begin{proof}
  \proves{propo:compact-cell-basis}
  \notready
We need to prove that for every ball \( B(x,r) \), the there is a cell
\( x\in\Gg_{s}\sbs B(x,r) \).
Let \( C \) be the complement of \( B(x,r) \).
For each point \( y \) of \( C \), there is an \( i \) such that
\( b_{i}(x)\cdot b_{i}(y) < 0 \).
In this case, let \( J^{0} = \setof{z: b_{i}(z) < 0} \),
\( J^{1} = \setof{z: b_{i}(z) > 0} \).
Let \( J(y)=J^{p} \) such that \( y\in J^{p} \).
Then \( C \sbs \bigcup_{y} J(y) \), and compactness implies that there is a
finite sequence \( y_{1},\dots,y_{k} \) with 
\( C \sbs \bigcup_{j=1}^{k} J(y_{j}) \).
Clearly, there is a cell 
 \[
  x \in \Gg_{s} \sbs B(x,r) \xcpt \bigcup_{j=1}^{k} J(y_{j}).
 \]
 \end{proof}


\subsection{Computable probability spaces}

If a separating sequence is given in advance, we may
restrict attention to the class of measures that make our sequence a
\( \mu \)-continuity sequence:

% Gacs source line 7209
 \begin{definition}
  \label{def:g4.7.8}
  \notready
Let \( \cM^{0}(X) \) be the set of those probability measures \( \mu \)
in \( \cM(X) \) for which \( \mu(X \xcpt X^{0})=0 \).
 \end{definition}

On the other hand, for each computable measure \( \mu \),
a computable separating sequence can be constructed that is a \( \mu \)-continuity
sequence.
Recall that \( B(x,r) \) is the ball of center \( x \) and radius \( r \).
Let \( D=\{s_{1},s_{2},\dots\}=\{\ag(1),\ag(2),\dots\} \) 
be the canonical enumeration of the canonical dense set \( D \).

% Gacs source line 7221
\begin{theorem}[Hoyrup-Rojas]\label{thm:continuity-balls}
  \notready
There is a sequence of uniformly computable reals \( (r_n)_{n\in \dN} \) such
that \( (B(s_{i},r_{n}))_{i,n} \) is a basis of balls that are \( \mu \)-continuity sets.
This basis is constructively equivalent to the original one, consisting of all 
balls \( B(s_{i},r) \), \( r\in\dQ \).
\end{theorem}

% Gacs source line 7228
 \begin{corollary}
  \label{cor:g4.7.9}
  \notready
There is a computable separating sequence with the \( \mu \)-continuity property.
 \end{corollary}
 \begin{proof}
  \proves{cor:g4.7.9}
  \notready
Let us list all balls \( B(s_{i},r_{n}) \) into a single sequence
\( B(s_{i_{k}},r_{n_{k}}) \).
The functions 
 \begin{align*}
  b_{k}(x)=d(s_{i_{k}},x)-r_{n_{k}}
 \end{align*}
give rise to the desired sequence.
 \end{proof}

For the proof of the theorem, we use some preparation.
Recall from Definition~A.2.5 that an atom is a point with positive
measure.

% Gacs source line 7245
\begin{lemma}\label{lem:zero_measure_points}
  \notready
Let \( X \) be \( \dR \) or \( \dR^+ \) or \( \clint{0}{1} \). 
Let \( \mu \) be a computable probability measure on \( X \). 
Then there is a sequence of uniformly computable reals \( (x_{n})_{n} \)
which is dense in \( X \) and contains no atoms of \( \mu \).
\end{lemma}
\begin{proof}
  \proves{lem:zero_measure_points}
  \notready
Let \( I \) be a closed rational interval. 
We construct \( x\in I \) with \( \mu(\{x\})=0 \). 
To do this, we construct inductively a nested sequence of
closed intervals \( J_k \) of measure \( <2^{-k+1} \), with \( J_0=I \). 
Suppose \( J_k=\clint{a}{b} \) has been constructed, with \( \mu(J_k)<2^{-k+1} \). 
Let
\( m=(b-a)/3 \): one of the intervals \( \clint{a}{a+m} \) and \( \clint{b-m}{b} \) must
have measure \( <2^{-k} \), and we can find it effectively---let it be \( J_{k+1} \).

From a constructive enumeration \( (I_n)_n \) of all the dyadic intervals, we
can construct \( x_n\in I_n \) uniformly.
 \end{proof}

% Gacs source line 7265
\begin{corollary}\label{coroll:sequence}
  \notready
Let \( (\bX,\mu) \) be a computable metric space with a computable measure
and let \( (f_i)_i \) be a sequence of uniformly computable real
valued functions on \( X \). 
Then there is a sequence of uniformly computable reals \( (x_n)_n \) that is
dense in \( \dR \) and such that each \( x_{n} \) is a \( \mu \)-continuity point of each
\( f_{i} \). 
\end{corollary}
\begin{proof}
  \proves{coroll:sequence}
  \notready
Consider the uniformly computable measures \( \mu_i=\mu\circ f_i^{-1} \) and define
\( \nu=\sum_i2^{-i}\mu_i \).
It is easy to see that \( \nu \) is a computable measure and then, by
Lemma~\ref{lem:zero_measure_points}, there is a sequence of uniformly computable
reals \( (x_n)_n \) which is dense in \( \dR \) and contains no atoms of \( \nu \). 
Since \( \nu(A)=0 \) iff \( \mu_i(A)=0 \) for all \( i \), we get \( \mu(\{f_i^{-1}(x_n)\})=0 \)
for all \( i,n \).
\end{proof}

\begin{proof}[Proof of Theorem~\protect\ref{thm:continuity-balls}]
  \proves{thm:continuity-balls}
  \notready
Apply Corollary~\ref{coroll:sequence} to \( f_i(x)=d(s_i,x) \).

Since every ball \( B(s_{i},r) \) can be expressed as a computably enumerable union of the balls
of the type \( B(s_{i},r_{n}) \) just constructed, the two bases are constructively
equivalent. 
\end{proof}
